%% ============================================================================
%% Chapter 9: API Versioning & Evolution Strategies
%% Mastering APIs: From Zero to Production Guru
%% ============================================================================
\chapter{API Versioning \& Evolution Strategies}
\label{ch:versioning}

%% ----------------------------------------------------------------------------
%% Executive Summary
%% ----------------------------------------------------------------------------
An API is a promise. The moment a single external consumer integrates against
your endpoints, every field name, every enum value, every error code, and every
default becomes a load-bearing part of someone else's production system --- and
you can no longer change any of it unilaterally without a plan. Yet APIs
\emph{must} change: businesses launch new products, regulations demand new
fields, performance work restructures aggregates, and yesterday's design
decisions become today's constraints. The central question of this chapter is
therefore not \emph{whether} to change an API, but \emph{how to change it
without breaking the consumers who trusted you}. Versioning is the discipline
that answers that question; evolution strategy is the engineering program that
makes the discipline affordable.

We begin by formalizing the vocabulary. A \emph{breaking change} is defined
precisely --- not by intuition but by the existence of a formerly valid client
interaction that the new contract rejects, mishandles, or renders
uninterpretable --- and we classify the everyday catalog of changes (removing
or renaming fields, tightening validation, changing types or semantics, adding
required fields, adding enum values, changing URLs) against that definition.
We examine the robustness principle (Postel's law) as a design posture for
APIs, and just as carefully examine its limits, because blind tolerance
fossilizes bugs into the contract. We then compare the four dominant
versioning mechanisms --- URI path, query parameter, custom header, and
\texttt{Accept} media-type negotiation --- with a rigorous trade-off matrix
covering caches, gateways, documentation, client pinning, and hypermedia, and
we study what Stripe, GitHub, and Google actually do in production and the
reasoning behind each choice~\cite{gehl2021apidesign}.

Mechanics follow policy. We study the expand-and-contract (parallel change)
pattern as the canonical zero-downtime migration shape, adapter and
anti-corruption layers that translate between version surfaces, dual-serving
architectures that run $n$ public versions over one service core, consumer
migration campaigns, and the usage telemetry that tells you when a version is
finally safe to retire. Deprecation is treated as an engineering discipline
with wire-level artifacts: the \texttt{Deprecation} and \texttt{Sunset}
response headers of RFC~8594, updated by RFC~9745~\cite{rfc8594},
\texttt{Link} relations that point at successor versions, sunset timelines,
contractual stability SLAs, SemVer discipline for API contracts, and
changelogs that consumers can actually plan against. Schema evolution gets its
own rulebook --- additive-only changes, never reusing names, enum
extensibility strategies, and the sharply different unknown-field tolerance
rules of JSON versus Protocol Buffers (Chapter~4)~\cite{protobufdocs} --- and
we close the loop with compatibility testing in CI: OpenAPI diffing in the
style of \texttt{oasdiff} as a breaking-change gate that fails the build
before a bad contract ever ships~\cite{openapi31}.

The code is deliberately production-shaped. You will build a dual-serving
FastAPI catalog API in which v1 is a pure adapter over the v2 domain model and
carries live \texttt{Deprecation}/\texttt{Sunset} headers; a stdlib-plus-YAML
OpenAPI diff gate that classifies changes as breaking or non-breaking and
exits non-zero on breakage; a media-type version negotiator with correct 406
behavior; a consumer-tolerance test suite that \emph{proves} old clients
survive additive change; and a retirement-readiness report generator that
turns access logs into a decommission decision. Everything is offline,
deterministic, typed Python~3.11, set in the fictional Northwind Robotics
parts-catalog universe.

\begin{keyconceptbox}
A \textbf{version} is not a number on a URL; it is a \emph{supported contract}
with a lifecycle: introduced, current, deprecated, sunset, retired. Breaking
changes demand a new version; non-breaking changes ride the existing one. The
cheapest version is the one you never had to mint --- prefer additive,
tolerant-reader evolution and reserve explicit versions for changes that are
genuinely incompatible. Retirement, not introduction, is the expensive part of
the lifecycle, so instrument usage from day one.
\end{keyconceptbox}

%% ----------------------------------------------------------------------------
\section{Learning Objectives \& Prerequisites}
\label{sec:ch9-objectives}

\subsection*{Learning Objectives}
After completing this chapter, its lab, and its exercises, you will be able to:

\begin{enumerate}[label=\textbf{LO-\arabic*.}, leftmargin=2.6cm]
  \item Classify any concrete API change --- field removal, rename, type
        change, semantic change, validation tightening, required-field
        addition, enum-value addition, error-code change, URL change --- as
        breaking or non-breaking, and defend the classification with the
        formal definition of a breaking change, including the subtle
        enum-addition caveat.
  \item State the robustness principle (Postel's law) as it applies to API
        design, apply tolerant-reader and additive-only disciplines, and
        articulate the principle's failure modes: bug fossilization, ambiguity
        amplification, and security surface growth.
  \item Compare URI path, query-parameter, custom-header, and
        \texttt{Accept} media-type versioning across caching, gateway routing,
        documentation, client pinning, hypermedia, and debuggability; explain
        why Stripe uses date-based versions, why GitHub versions via a header,
        and why Google versions in the URI~\cite{gehl2021apidesign}.
  \item Execute an expand-and-contract migration end to end: deploy the new
        contract alongside the old, dual-serve through an adapter layer,
        migrate consumers with telemetry-driven campaigns, and contract by
        retiring the old version only when the data permits it.
  \item Engineer a deprecation policy: emit RFC~8594/RFC~9745
        \texttt{Deprecation} and \texttt{Sunset} headers, publish sunset
        timelines and stability SLAs, maintain a SemVer-disciplined changelog,
        and select communication channels proportionate to blast radius
        \cite{rfc8594}.
  \item Apply schema-evolution rules --- additive-only changes, never reusing
        names, enum extensibility with fallback values, unknown-field
        tolerance --- and contrast JSON's name-based tolerance with Protocol
        Buffers' tag-based rules~\cite{protobufdocs}.
  \item Build, in typed Python~3.11, a dual-serving FastAPI application with a
        v1-to-v2 adapter layer, an OpenAPI breaking-change CI gate, a
        media-type version negotiator, a consumer-tolerance test suite, and a
        retirement-readiness report generator~\cite{fastapidocs,openapi31}.
  \item Design a version-support policy at scale: decide how many versions to
        run concurrently, quantify the per-version carrying cost, and
        differentiate policy for internal versus external APIs.
\end{enumerate}

\subsection*{Prerequisites}
This chapter assumes you have completed:

\begin{itemize}
  \item \textbf{Chapter 0} --- environment setup: Python 3.11+, a virtual
        environment, and the ability to run \texttt{pip} and \texttt{pytest}
        from PowerShell.
  \item \textbf{Chapter 1} --- the HTTP wire protocol: methods, status codes,
        and header fields, especially \texttt{Accept}, \texttt{Content-Type},
        \texttt{Link}, and cache-related headers~\cite{rfc9110}.
  \item \textbf{Chapter 2} --- consuming APIs from Python: parsing JSON
        responses and handling HTTP errors.
  \item \textbf{Chapter 3} --- the REST architectural style: resources,
        representations, and hypermedia, which inform the
        URI-versus-media-type debate.
  \item \textbf{Chapter 4} --- serialization formats and data contracts:
        JSON Schema, Protocol Buffers tag rules, content negotiation, and the
        backward/forward/full compatibility vocabulary we extend here
        \cite{protobufdocs}.
  \item \textbf{Chapter 5} --- client-side security fundamentals, and
        \textbf{Chapter 6} --- building APIs with FastAPI: routing, Pydantic
        models, and \texttt{TestClient}~\cite{fastapidocs}.
\end{itemize}

No network access, cloud account, or API key is required anywhere in this
chapter: every fixture is synthetic, embedded, and deterministic, set in the
fictional \emph{Northwind Robotics} parts and logistics universe.

\begin{warningbox}
Versioning policy is a \textbf{business commitment} encoded in engineering
artifacts. A \texttt{Sunset} header, a changelog entry, and a support SLA are
promises your consumers will plan around. Do not publish a deprecation date
you are not prepared to honor, and do not retire a version while consumers
you have not contacted still depend on it. The code in this chapter is simple;
the discipline around it is the actual subject.
\end{warningbox}

%% ----------------------------------------------------------------------------
\section{Deep Technical Mechanics \& Architectural Concepts}
\label{sec:ch9-mechanics}

\subsection{A Taxonomy of Change}
\label{subsec:ch9-taxonomy}

Every versioning decision reduces to one question: \emph{does this change
break an existing consumer?} Answering it honestly requires a definition
precise enough to classify borderline cases, not a vibe.

\begin{definition}[Breaking change]
\label{def:breaking}
Let $C$ be an API contract and $\mathcal{K}(C)$ the set of client interactions
(request plus response interpretation) that are valid and succeed under $C$.
A change producing contract $C'$ is \emph{breaking} if and only if there
exists at least one interaction $k \in \mathcal{K}(C)$ such that under $C'$,
$k$ is rejected, fails, or returns a response that a client written against
$C$ cannot interpret with its existing parsing and handling logic. A change
with no such $k$ is \emph{non-breaking} (compatible).
\end{definition}

\begin{definition}[Compatibility modes]
\label{def:compat}
A new contract (or schema) version $C'$ is \emph{backward compatible} with
$C$ if consumers written against $C$ work unchanged against $C'$ (new server,
old clients). $C'$ is \emph{forward compatible} with $C$ if consumers written
against $C'$ can meaningfully process data produced under $C$ (old readers,
new data). A change that is both is \emph{fully compatible}. Web API evolution
is overwhelmingly concerned with backward compatibility; durable data formats
(Chapter~4) must care about both directions.
\end{definition}

With Definition~\ref{def:breaking} in hand, the everyday catalog of changes
classifies cleanly --- with two notorious edge cases that we flag explicitly.

\begin{table}[ht]
  \centering
  \small
  \begin{tabularx}{\textwidth}{l l X}
    \toprule
    \textbf{Change} & \textbf{Verdict} & \textbf{Why} \\
    \midrule
    Remove a response field & Breaking & Clients reading the field fail or
      misinterpret; violates the old contract's response shape. \\
    Rename a field & Breaking & A rename is a removal plus an addition; the
      removal half breaks. ``Harmless renames'' are the most common
      accidental breakage. \\
    Change a field's type & Breaking & \texttt{price} as float dollars
      becoming integer cents deserializes wrongly or fails validation on
      every old client. \\
    Change field semantics & Breaking & Same name, same type, new meaning
      (e.g., \texttt{stock} switches from on-hand to available-to-promise):
      old clients silently compute wrong answers. Worst class of break ---
      no error is raised. \\
    Tighten validation & Breaking & Requests that succeeded now return 400.
      Loosening validation is non-breaking. \\
    Add a required request field & Breaking & Old clients cannot supply it;
      their formerly valid requests are rejected. \\
    Add an optional request field & Non-breaking & Old requests remain valid;
      the server defaults the new field. \\
    Add a response field & Non-breaking\textsuperscript{*} & Tolerant readers
      ignore it; strict-schema clients may not --- see below. \\
    Add an endpoint & Non-breaking & Nothing existing changed; new capability
      only. \\
    Change a URL path & Breaking & Every hard-coded client reference 404s.
      URL \emph{additions} are fine; URL \emph{moves} break. \\
    Change an error code/body & Breaking & Clients branch on status codes and
      problem-detail members; changing them breaks error-handling paths. \\
    Add an enum value & \textbf{Conditional} & Breaking for clients that
      validate exhaustively or \texttt{switch} without a default arm;
      non-breaking for tolerant readers. See the caveat below. \\
    \bottomrule
  \end{tabularx}
  \caption{Change taxonomy under Definition~\ref{def:breaking}.}
  \label{tab:ch9-taxonomy}
\end{table}

\paragraph{The enum-addition caveat.}
Adding a value to an enumeration is the most underestimated breaking change
in API design. Producers think of enums as closed sets they own; consumers
frequently compile them into closed \texttt{switch}/\texttt{match} statements,
strict deserializers, or database check constraints. When Northwind Robotics
adds \texttt{battery} to the catalog \texttt{category} enum, every mobile
client whose generated model rejects unknown enum values does not degrade ---
it \emph{crashes}, typically inside its JSON decoding layer, on an otherwise
perfectly valid response. The mitigation is contractual: document enums as
\emph{open} (extensible) from day one, require clients to define fallback
behavior for unknown values, and provide a sentinel such as
\texttt{unknown}/\texttt{other} in the published schema. Protocol Buffers
codifies exactly this posture: unrecognized enum values are preserved and
surfaced as unknown values rather than rejected~\cite{protobufdocs}. If your
existing consumer base already treats the enum as closed, adding a value
\emph{is} a breaking change and belongs in a new version.

\subsection{The Robustness Principle and Its Limits}
\label{subsec:ch9-postel}

Postel's law --- \emph{be conservative in what you send, be liberal in what
you accept} --- translates into two concrete API disciplines. On the producer
side, \textbf{conservative sending} means additive-only evolution, stable
field names, documented-open enums, and never reusing a retired name for a
new meaning. On the consumer side, \textbf{liberal acceptance} means the
\emph{tolerant reader}: ignore fields you do not recognize, do not fail on
unexpected enum values, do not depend on object member ordering, and treat
absence of an optional field as a normal case. A consumer base built on
tolerant readers makes most additive producer changes free, which is why the
tolerant-reader posture is the single highest-leverage compatibility
investment an ecosystem can make.

The principle has real limits, and a principal engineer must know them.
First, \textbf{bug fossilization}: every quirk you silently tolerate ---
the misspelled field you alias, the malformed date you repair --- becomes a
de facto contract you can never remove, because some consumer now depends on
your tolerance. Second, \textbf{ambiguity amplification}: accepting
\texttt{price} as either a string or a number forces every consumer to guess
which semantics apply, multiplying interoperability bugs rather than
preventing them. Third, \textbf{security surface}: liberal parsers are
historically fertile ground for injection, confusion, and denial-of-service
attacks; validate strictly at the trust boundary (Chapter~5) and be liberal
only about \emph{ignoring} the unrecognized, never about \emph{accepting} the
malformed. The mature reading of Postel is: strict about structure and safety,
tolerant about extensibility.

\subsection{Versioning Strategies: A Rigorous Comparison}
\label{subsec:ch9-strategies}

When a change is genuinely breaking, you must mint a new version, and the
version identifier must live somewhere on the wire. Four placements dominate
practice; Figure~\ref{fig:ch9-routing} shows the selection flow they share.

\paragraph{URI path versioning.} The version is a path segment:
\texttt{/v1/catalog/WR-1042}. Strengths: maximally visible and debuggable
(the version is obvious in every log line and browser address bar); routing
is trivial for any gateway; per-version cache keys emerge naturally because
the URI is the primary cache key in HTTP~\cite{rfc9110}; client pinning is
effortless. Weaknesses: it violates the REST ideal that a URI identifies a
\emph{resource}, not a representation of it --- \texttt{/v1/catalog/WR-1042}
and \texttt{/v2/catalog/WR-1042} are the same part wearing two addresses;
links embedded in stored data (bookmarks, database columns, emails) pin
consumers to versions; hypermedia payloads must be regenerated per version.

\paragraph{Query-parameter versioning.} \texttt{/catalog/WR-1042?version=2}.
Largely the URI strategy with worse ergonomics: caches \emph{do} key on query
strings, so cacheability survives, but intermediaries and logging pipelines
more often normalize or drop query parameters than path segments, and the
version reads as optional decoration rather than structural commitment.

\paragraph{Custom header versioning.} \texttt{X-API-Version: 2}. Keeps URIs
stable and pure, but suffers at the edges: browser debugging is awkward,
many gateways require explicit configuration to route on custom headers,
documentation tooling rarely models it, and --- critically --- naive caches
keyed only on URI will serve a v1 body to a v2 request unless the server
emits \texttt{Vary: X-API-Version} and every intermediary honors it
\cite{rfc9110}. Cache poisoning across versions is a real, observed failure
mode (Section~\ref{sec:ch9-pitfalls}).

\paragraph{Media-type (\texttt{Accept} header) versioning.}
\texttt{Accept: application/vnd.northwind.catalog+json; version=2}. This is
the most REST-orthodox placement: the URI identifies the resource; the
version is a property of the negotiated \emph{representation}, exactly as
content negotiation intends (Chapter~4)~\cite{rfc9110}. It composes naturally
with hypermedia and permits resource-by-resource versioning granularity. Its
costs are the sharpest: the highest complexity floor for consumers,
near-invisibility in logs and browsers, gateway routing that must parse
media-type parameters, the same \texttt{Vary: Accept} cache requirement, and
documentation tooling that assumes path-based APIs. Listing~9.3 implements a
correct negotiator, including the 406 response for unsupported versions.

\begin{table}[ht]
  \centering
  \small
  \begin{tabularx}{\textwidth}{X c c c c}
    \toprule
    \textbf{Criterion} & \textbf{URI path} & \textbf{Query param} &
      \textbf{Custom header} & \textbf{Media type} \\
    \midrule
    Cache correctness & Excellent & Good & Fragile (needs \texttt{Vary}) &
      Fragile (needs \texttt{Vary}) \\
    Gateway/LB routing & Trivial & Easy & Config required & Config required \\
    Debuggability (browser, logs) & Excellent & Good & Poor & Poor \\
    Client pinning difficulty & Trivial & Easy & Moderate & Hard \\
    Hypermedia fit & Poor & Poor & Good & Excellent \\
    Doc/tooling support & Excellent & Good & Fair & Fair \\
    REST resource purity & Poor & Fair & Good & Excellent \\
    Version granularity & Per API & Per API & Per API & Per resource \\
    \bottomrule
  \end{tabularx}
  \caption{Versioning strategy trade-off matrix.}
  \label{tab:ch9-matrix}
\end{table}

\begin{figure}[ht]
\centering
\begin{tikzpicture}[
  font=\small, >=stealth, node distance=0.55cm and 1.1cm,
  box/.style={rectangle, draw, rounded corners, align=center,
    minimum width=3.4cm, minimum height=0.9cm},
  decision/.style={diamond, draw, align=center, aspect=2.2, inner sep=1pt},
  term/.style={rectangle, draw, fill=gray!15, rounded corners, align=center,
    minimum width=2.6cm, minimum height=0.8cm}]
  \node[box] (req) {HTTP request};
  \node[box, right=of req] (extract) {Extract version signal\\ \footnotesize URI segment / query / header / \texttt{Accept}};
  \node[decision, right=of extract] (known) {Signal\\ recognized?};
  \node[decision, below=of known] (supported) {Version\\ supported?};
  \node[term, right=of supported] (reject) {406 Not Acceptable\\ \footnotesize (+\ \texttt{Link} to supported versions)};
  \node[box, left=of supported] (route) {Route to version handler\\ \footnotesize direct (v2) or via adapter (v1)};
  \node[term, below=of route] (serve) {Serve from shared service core};
  \draw[->] (req) -- (extract);
  \draw[->] (extract) -- (known);
  \draw[->] (known) -- node[right] {yes} (supported);
  \draw[->] (known.east) -| node[pos=0.25, above] {no} ($(reject.north)+(0,0.55)$) -- (reject.north);
  \draw[->] (supported) -- node[above] {yes} (route);
  \draw[->] (supported) -- node[above] {no} (reject);
  \draw[->] (route) -- (serve);
\end{tikzpicture}
\caption{Version-selection routing flow. An unrecognized \emph{signal} and an
unsupported \emph{version} are different failures; both surface as 406 (or 404
for unknown paths) with pointers to supported versions.}
\label{fig:ch9-routing}
\end{figure}

\paragraph{What major providers do, and why.}
\textbf{Stripe} versions by \emph{date}: every account is pinned to a version
like \texttt{2024-06-20}, sent per request via the \texttt{Stripe-Version}
header, and each response is rendered through a cascade of
backward-compatibility transformation layers from current internals down to
the caller's pinned date. The payoff is that a consumer's integration
\emph{never changes underneath them} unless they explicitly upgrade, while
Stripe ships continuous internal change. \textbf{GitHub} versions via a media
type plus an \texttt{X-GitHub-Api-Version} date header, keeping resource URIs
stable across versions --- the REST-orthodox choice, suited to a
hypermedia-rich API. \textbf{Google} versions in the URI (\texttt{/v1/},
\texttt{/v2/}) and codifies the policy in its API Improvement Proposals:
discoverability, gateway simplicity, and cache clarity outweigh URI purity at
Google's operational scale~\cite{gehl2021apidesign}. The pattern to
internalize is not \emph{which} mechanism won but \emph{why}: each provider
chose the placement whose costs its consumer base and infrastructure could
bear, and each pairs the mechanism with a strict non-breaking-change policy
so that new versions are rare events.

\subsection{Evolution Mechanics: Expand-and-Contract and Dual-Serving}
\label{subsec:ch9-expandcontract}

The \textbf{expand-and-contract} pattern (also called \emph{parallel change})
is the canonical zero-downtime migration shape for any shared contract, from
database schemas~\cite{kleppmann2017ddia} to public APIs. Its discipline is
that the system passes through an intermediate state in which \emph{both} the
old and new contracts are simultaneously valid, so no consumer ever observes
a discontinuity. The pattern has three phases, shown in
Figure~\ref{fig:ch9-timeline}:

\begin{enumerate}
  \item \textbf{Expand.} Add the new contract (endpoint version, field, or
        representation) alongside the old one. Nothing is removed. Old
        consumers are unaffected by construction; new consumers can opt in.
  \item \textbf{Migrate.} Move consumers from old to new, at their pace but
        on your schedule: announce deprecation, publish a migration guide,
        instrument usage, and run a migration campaign
        (Section~\ref{subsec:ch9-campaigns}).
  \item \textbf{Contract.} When telemetry shows the old surface is unused (or
        below a contractual threshold with stragglers individually engaged),
        retire it. Only now is the old code deleted.
\end{enumerate}

\begin{figure}[ht]
\centering
\begin{tikzpicture}[font=\small, >=stealth]
  % Timeline axis
  \draw[->, thick] (0,0) -- (14.2,0) node[right] {time};
  % Phase boundary ticks
  \foreach \x/\lab in {0/{$T_0$}, 4.2/{$T_1$}, 9.4/{$T_2$}, 13.6/{$T_3$}} {
    \draw (\x, 0.12) -- (\x, -0.12);
    \node[below] at (\x, -0.15) {\lab};
  }
  % Phase labels above the axis
  \node[align=center] at (2.1, 0.85) {\textbf{Expand}\\ \footnotesize deploy v2 beside v1};
  \node[align=center] at (6.8, 0.85) {\textbf{Migrate}\\ \footnotesize consumers move; v1 deprecated};
  \node[align=center] at (11.5, 0.85) {\textbf{Contract}\\ \footnotesize retire v1};
  % Serving-state braces below
  \draw[decorate, decoration={brace, amplitude=5pt, mirror}]
    (0.1, -1.0) -- (4.1, -1.0) node[midway, below=6pt, note] {v1 only};
  \draw[decorate, decoration={brace, amplitude=5pt, mirror}]
    (4.3, -1.0) -- (13.5, -1.0) node[midway, below=6pt, note] {dual-serving: v1 (deprecated) + v2};
  \draw[decorate, decoration={brace, amplitude=5pt}]
    (13.7, 0.45) -- (14.1, 0.45) node[midway, above=6pt, note] {v2 only};
  % Key events
  \node[note, above] at (0.0, 1.75) {v1 stable};
  \node[note, above, align=center] at (4.2, 1.75) {v2 GA;\\ \texttt{Deprecation} on v1};
  \node[note, above, align=center] at (9.4, 1.75) {v1 usage $<$ threshold;\\ sunset countdown};
  \node[note, above, align=center] at (13.6, 1.75) {\texttt{Sunset} date;\\ v1 returns 410};
\end{tikzpicture}
\caption{Expand-and-contract migration timeline. The dual-serving window
between $T_1$ and $T_3$ is where all consumer risk is absorbed; contraction
before the window closes is what causes outages.}
\label{fig:ch9-timeline}
\end{figure}

\paragraph{Adapters and anti-corruption layers.}
Dual-serving does \emph{not} mean maintaining two implementations of the
business logic. The architecture that keeps it affordable is one service core
expressed in the \emph{newest} domain model, with each older public version
implemented as a thin \textbf{adapter} --- a pure translation layer mapping
between the old wire contract and the current model. The adapter is an
anti-corruption layer in the domain-driven-design sense: it prevents legacy
vocabulary (\texttt{price} in float dollars, \texttt{in\_stock} as a boolean)
from leaking into the current model (\texttt{unit\_price\_cents},
\texttt{stock} count). Adapters must be pure, total, and tested as contracts:
every field in the old surface must be derivable from the new model, and the
test suite (Listing~9.4) must prove old clients cannot observe the
translation. When a field \emph{cannot} be faithfully derived --- say the v2
model drops a concept v1 exposed --- that is the signal that retirement of
v1, not another adapter hack, is the honest path. Figure~\ref{fig:ch9-dual}
shows the serving architecture; Listing~9.1 implements it.

\begin{figure}[ht]
\centering
\begin{tikzpicture}[
  font=\small, >=stealth, node distance=0.6cm and 1.0cm,
  client/.style={rectangle, draw, rounded corners, align=center,
    minimum width=2.4cm, minimum height=0.8cm},
  layer/.style={rectangle, draw, align=center, minimum width=4.6cm,
    minimum height=0.95cm},
  corebox/.style={rectangle, draw, align=center, minimum width=4.6cm,
    minimum height=1.1cm, fill=gray!12}]
  % Clients
  \node[client] (web) {Legacy mobile app\\ \footnotesize (v1)};
  \node[client, below=of web] (partner) {Partner integration\\ \footnotesize (v1)};
  \node[client, below=of partner] (modern) {New web storefront\\ \footnotesize (v2)};
  % Gateway
  \node[layer, right=1.6cm of partner] (gw) {API gateway / router\\ \footnotesize version selection (Figure~\ref{fig:ch9-routing})};
  % Version surfaces
  \node[layer, right=1.6cm of gw, yshift=1.1cm] (v1) {v1 adapter\\ \footnotesize translates to v1 wire model\\ \footnotesize +\ \texttt{Deprecation}/\texttt{Sunset} headers};
  \node[layer, below=0.6cm of v1] (v2) {v2 controller\\ \footnotesize native wire model};
  % Core
  \node[corebox, below=0.9cm of v2] (core) {Service core (v2 domain model)\\ \footnotesize catalog logic, validation, pricing};
  \node[corebox, below=0.5cm of core, minimum height=0.8cm] (db) {Datastore};
  \node[draw, dashed, inner sep=0.35cm, fit=(v1)(v2)(core)(db),
    label={[note]above right:single deployable --- one core, $n$ surfaces}] {};
  % Edges
  \draw[->] (web.east) -- ++(0.5,0) |- (gw.west);
  \draw[->] (partner) -- (gw);
  \draw[->] (modern.east) -- ++(0.5,0) |- (gw.west);
  \draw[->] (gw.east) |- (v1.west);
  \draw[->] (gw.east) |- (v2.west);
  \draw[->] (v1) -- (core);
  \draw[->] (v2) -- (core);
  \draw[<->] (core) -- (db);
  % Telemetry tap
  \node[client, below=1.5cm of gw, minimum width=3.4cm] (tele) {Usage telemetry\\ \footnotesize per-version counts $\to$ Listing 9.5};
  \draw[->, dashed] (gw) -- (tele);
\end{tikzpicture}
\caption{Dual-serving architecture. Both public versions are surfaces over one
service core; v1 is a pure adapter with deprecation metadata, and the gateway
feeds per-version telemetry that later justifies retirement.}
\label{fig:ch9-dual}
\end{figure}

\paragraph{Consumer migration campaigns.}
\label{subsec:ch9-campaigns}
A deprecation announcement is not a migration. A migration campaign is a
managed program with an owner, a consumer registry, and a deadline. Identify
consumers from telemetry (API keys, OAuth client IDs, IP ranges as a last
resort); segment by traffic and criticality; contact the long tail
individually --- an email to a consumer doing 40\%\ of your v1 traffic is
worth more than a banner seen by everyone; publish a migration guide with a
field-level mapping table and copy-pasteable diffs; offer a feedback channel
and a published escalation path; and schedule brownouts (short, announced
windows in which the old version returns errors) as a forcing function for
consumers who do not read email. Brownouts convert silent dependency into
visible, survivable failure while there is still time to act.

\paragraph{Usage telemetry as the retirement oracle.}
The contract phase can only begin when you can \emph{prove} who still depends
on the old version. Log, at minimum: version, endpoint, consumer identity,
and timestamp on every request. From that stream compute per-version request
share, distinct-consumer counts, and trend. Listing~9.5 turns such a log into
a retirement-readiness verdict. A common threshold policy: schedule sunset
when the old version drops below 5\%\ of traffic \emph{and} every remaining
consumer has been contacted; execute sunset only when remaining consumers
have migrated or formally accepted the cutoff.

\subsection{Deprecation Policy Engineering}
\label{subsec:ch9-deprecation}

Deprecation is a state, not a deletion. A deprecated version still works; it
is no longer recommended, receives no enhancements, and has a published
end-of-life. The wire-level vocabulary is standardized. RFC~8594 introduced
the \texttt{Deprecation} and \texttt{Sunset} response header fields
\cite{rfc8594}; RFC~9745 (2025) subsequently revised \texttt{Deprecation},
replacing the original HTTP-date value with a Unix-timestamp form
\texttt{@}\emph{<seconds>}, while \texttt{Sunset} retains an HTTP-date. A
well-behaved deprecated endpoint therefore answers like Listing~9.1's v1
surface:

\begin{itemize}
  \item \texttt{Deprecation: @1788220800} --- the moment deprecation began
        (RFC~9745 form; RFC~8594-era deployments used an HTTP-date here).
  \item \texttt{Sunset: Sat, 01 May 2027 00:00:00 GMT} --- the committed
        moment the version becomes unresponsive.
  \item \texttt{Link: <https://api.example/v2/catalog>;
        rel="successor-version"} and
        \texttt{rel="deprecation"} pointing at the migration guide ---
        machine-followable pointers, not just prose.
\end{itemize}

Around those headers sits the policy. \textbf{Sunset timelines} must be
generous and public: twelve months is a common floor for external APIs with
managed consumers; six may suffice for a fast-moving partner API; internal
APIs can contract in weeks. \textbf{Stability SLAs} put the promise in the
contract: ``no breaking changes within a major version,'' ``deprecated
versions receive security fixes until sunset,'' ``at least 12 months between
deprecation and sunset.'' \textbf{SemVer for APIs} gives the changelog its
grammar: MAJOR for breaking changes (a new version), MINOR for
backward-compatible additions (new field, new endpoint, new \emph{documented
open} enum value), PATCH for compatible fixes. \textbf{Changelog discipline}
means every change is dated, classified by that grammar, and written for the
consumer (\emph{``field X added to resource Y''}), not the implementer.
\textbf{Communication channels} scale with blast radius: headers and
changelogs for the attentive, developer-portal banners and newsletters for
the reachable, direct contact for the identified, brownouts for the rest.

\subsection{Schema Evolution Rules}
\label{subsec:ch9-schema}

Within a version, the schema may only move in compatible directions. The
rulebook:

\begin{enumerate}
  \item \textbf{Additive only.} New fields, new endpoints, new optional
        parameters. Never remove, never rename, never retype within a
        version.
  \item \textbf{Never reuse names.} A retired name is reserved forever;
        reintroducing \texttt{price} with new semantics corrupts every
        stored payload and cached response that still carries the old
        meaning. Protocol Buffers encodes this as \texttt{reserved} field
        numbers and names --- adopt the same discipline in JSON contracts
        even though the format cannot enforce it~\cite{protobufdocs}.
  \item \textbf{Enums are open by contract.} Publish them as extensible,
        ship a sentinel value, and require client fallback behavior. Adding
        a value to an enum documented as closed is a breaking change.
  \item \textbf{Unknown fields are tolerated, not errors.} This is where the
        wire format matters. In JSON, tolerance is a \emph{client
        discipline}: nothing in the format stops a strict deserializer from
        rejecting unknown members, so the contract must ask clients to
        ignore them (and generated models must be configured to do so). In
        Protocol Buffers, tolerance is a \emph{format property}: unknown
        tags are skipped by the wire parser itself, and since proto3 they
        are even retained on reserialize, so a proxy can pass through fields
        it does not understand~\cite{protobufdocs}. JSON gives you
        tolerance by convention; Protobuf gives it to you by construction
        (Chapter~4).
\end{enumerate}

\paragraph{Compatibility testing in CI.}
Rules without enforcement are wishes. The enforcement mechanism is a
breaking-change gate in the pipeline: diff the proposed OpenAPI description
against the published one, classify every delta, and fail the build on any
breaking change that is not accompanied by a version bump~\cite{openapi31}.
This is what tools such as \texttt{oasdiff} automate; Listing~9.2 builds a
minimal but real gate with stdlib plus PyYAML so the classification logic is
fully inspectable. The gate belongs on the pull request, where the cost of
the conversation is lowest --- discovering a break in code review costs
minutes; discovering it from a consumer's pager costs trust.

\subsection{Version Policy at Scale}
\label{subsec:ch9-scale}

Every concurrently supported version carries a recurring cost: adapter code
and its tests, a documentation surface, a compatibility CI lane, on-call
knowledge, and the combinatorial test matrix against every other moving part
of the system. A useful mental model:

\begin{equation}
C_{\text{carry}}(n) \;=\; n \cdot c_{\text{adapter}} \;+\; c_{\text{test}} \cdot \binom{n}{2} \;+\; c_{\text{docs}} \cdot n \;+\; c_{\text{incident}}(n)
\label{eq:ch9-cost}
\end{equation}

\noindent where $n$ is the number of live versions. The quadratic term ---
pairwise interaction testing --- is why ``just keep every version forever''
collapses operationally long before it collapses financially. Practical
policies: \textbf{two live versions} (current plus previous) for most
external APIs; \textbf{three} for platforms with slow-moving enterprise
consumers, with the third on an accelerated sunset; date-based schemes like
Stripe's bound $n$ by time rather than count (``all versions younger than 24
months''), which caps cost automatically. \textbf{Internal APIs} may adopt
much shorter windows --- weeks to months --- because consumers are
enumerable, reachable, and often deployable by the same organization;
coordinated expand-and-contract with a two-week dual-serving window is
routine. \textbf{External APIs} must assume unenumerable consumers, cached
copies of old documentation, and contracts you cannot see; timelines stretch
to quarters or years, and retirement requires the full campaign machinery of
Section~\ref{subsec:ch9-campaigns}. The universal rule: the number of
versions you support is a product decision with an engineering budget, not
an engineering convenience.

%% ----------------------------------------------------------------------------
\section{Production-Ready Code Implementation}
\label{sec:ch9-code}

All listings are typed Python~3.11, offline, deterministic, and set in the
Northwind Robotics catalog universe. Install the dependencies first
(Section~\ref{sec:ch9-deps}), place each listing in a working directory such
as \texttt{C:\textbackslash Training\textbackslash ch09}, and run from
PowerShell.

\subsection{Listing 9.1 --- Dual-Serving Catalog API (v1 + v2 over One Core)}
\label{subsec:ch9-listing1}

One service core holds the v2 domain model; the v1 surface is a pure adapter
with live \texttt{Deprecation}, \texttt{Sunset}, and \texttt{Link} headers.
Running the module executes a deterministic demonstration against FastAPI's
\texttt{TestClient} --- no server, no network~\cite{fastapidocs}.

\begin{minted}{python}
"""Dual-serving Catalog API (v1 + v2) for Northwind Robotics.

Chapter 9, Listing 9.1 -- one v2 domain model serves two public versions.
The v1 surface is a pure adapter over the v2 model and carries RFC 8594 /
RFC 9745 deprecation metadata. Offline: run the module for a deterministic
demonstration via FastAPI's TestClient.
"""

from __future__ import annotations

import enum
from datetime import datetime, timezone

from fastapi import FastAPI, Request, Response
from fastapi.responses import JSONResponse
from fastapi.testclient import TestClient
from pydantic import BaseModel, Field

# --- Deprecation policy constants (dates are fixed => deterministic output) ---
DEPRECATED_AT = datetime(2026, 9, 1, tzinfo=timezone.utc)
SUNSET_AT = datetime(2027, 5, 1, tzinfo=timezone.utc)
DEPRECATION_HEADER = f"@{int(DEPRECATED_AT.timestamp())}"          # RFC 9745
SUNSET_HEADER = "Sat, 01 May 2027 00:00:00 GMT"                    # RFC 8594
MIGRATION_GUIDE = "https://docs.northwind-robotics.example/catalog/migration"


# --- Single source of truth: the v2 domain model ------------------------------
class Category(str, enum.Enum):
    """Open enum: consumers MUST tolerate unknown values (fallback rule)."""

    ACTUATOR = "actuator"
    SENSOR = "sensor"
    CONTROLLER = "controller"
    BATTERY = "battery"  # added mid-lifecycle; safe only because enum is open


class PartV2(BaseModel):
    """Current (v2) domain model. All business logic speaks only this type."""

    part_id: str
    sku: str
    name: str
    unit_price_cents: int = Field(ge=0)
    currency: str = "USD"
    category: Category
    stock: int = Field(ge=0)
    tags: list[str] = []
    status: str = "active"


class PartV1(BaseModel):
    """Legacy (v1) wire contract, frozen forever: float-dollar price, bool stock."""

    id: str
    sku: str
    name: str
    price: float
    category: str
    in_stock: bool


CATALOG: dict[str, PartV2] = {
    p.part_id: p
    for p in [
        PartV2(part_id="WR-1042", sku="SRV-T450", name="Torque Servo T450",
               unit_price_cents=12999, category=Category.ACTUATOR, stock=34,
               tags=["servo", "high-torque"]),
        PartV2(part_id="WR-2210", sku="LDR-M90", name="Lidar Module M90",
               unit_price_cents=45900, category=Category.SENSOR, stock=0,
               tags=["lidar", "navigation"]),
        PartV2(part_id="WR-3307", sku="CTL-AX12", name="Axis Controller AX12",
               unit_price_cents=8900, category=Category.CONTROLLER, stock=112,
               tags=["controller", "can-bus"]),
        PartV2(part_id="WR-4100", sku="BAT-LP38", name="LiPo Pack LP38",
               unit_price_cents=15450, category=Category.BATTERY, stock=57,
               tags=["battery", "38wh"]),
    ]
}


# --- Anti-corruption layer: v2 domain model -> frozen v1 wire contract --------
def to_v1(part: PartV2) -> PartV1:
    """Pure, total translation: every v1 field is derivable from the v2 model."""

    return PartV1(
        id=part.part_id,
        sku=part.sku,
        name=part.name,
        price=round(part.unit_price_cents / 100, 2),
        category=part.category.value,
        in_stock=part.stock > 0,
    )


def lookup(part_id: str) -> PartV2 | None:
    return CATALOG.get(part_id.upper())


app = FastAPI(title="Northwind Robotics Catalog API", version="2.0.0")


@app.middleware("http")
async def deprecation_middleware(request: Request, call_next) -> Response:
    """Stamp RFC 8594/9745 lifecycle metadata on every v1 response."""

    response: Response = await call_next(request)
    if request.url.path.startswith("/v1/"):
        response.headers["Deprecation"] = DEPRECATION_HEADER
        response.headers["Sunset"] = SUNSET_HEADER
        response.headers["Link"] = (
            f'<{MIGRATION_GUIDE}>; rel="deprecation", '
            f'</v2/catalog>; rel="successor-version"'
        )
    return response


def not_found(version: str, part_id: str) -> JSONResponse:
    return JSONResponse(
        status_code=404,
        content={"type": "about:blank", "title": "Part not found",
                 "status": 404, "detail": f"no part '{part_id}' in {version}"},
    )


# --- v1 surface (deprecated): adapter over the core ---------------------------
@app.get("/v1/catalog", response_model=list[PartV1])
def list_catalog_v1() -> list[PartV1]:
    return [to_v1(p) for p in CATALOG.values()]


@app.get("/v1/catalog/{part_id}", response_model=PartV1)
def get_part_v1(part_id: str) -> PartV1 | JSONResponse:
    part = lookup(part_id)
    return to_v1(part) if part else not_found("v1", part_id)


# --- v2 surface (current): native domain model --------------------------------
@app.get("/v2/catalog", response_model=list[PartV2])
def list_catalog_v2() -> list[PartV2]:
    return list(CATALOG.values())


@app.get("/v2/catalog/{part_id}", response_model=PartV2)
def get_part_v2(part_id: str) -> PartV2 | JSONResponse:
    part = lookup(part_id)
    return part if part else not_found("v2", part_id)


def demo() -> None:
    client = TestClient(app)
    r1 = client.get("/v1/catalog/WR-1042")
    print("== v1 (deprecated) ==")
    print("Deprecation:", r1.headers.get("Deprecation"))
    print("Sunset:     ", r1.headers.get("Sunset"))
    print("Link:       ", r1.headers.get("Link"))
    print("Body:       ", r1.json())
    r2 = client.get("/v2/catalog/WR-1042")
    print("\n== v2 (current) ==")
    print("Deprecation:", r2.headers.get("Deprecation"))
    print("Body:       ", r2.json())
    r3 = client.get("/v2/catalog/NOPE-1")
    print("\n== unknown part ==", r3.status_code, r3.json()["detail"])


if __name__ == "__main__":
    demo()
\end{minted}

\subsection{Listing 9.2 --- OpenAPI Breaking-Change Gate (CI in a Box)}
\label{subsec:ch9-listing2}

The gate diffs two OpenAPI~3.x YAML descriptions, classifies every delta
against the taxonomy of Table~\ref{tab:ch9-taxonomy}, prints a report, and
exits non-zero on any breaking change --- the exact behavior you wire into a
pull-request pipeline~\cite{openapi31}. It deliberately handles the
high-signal subset (paths, operations, request bodies, response schemas) and
assumes inline schemas; resolving \texttt{\$ref} is Capstone~3.

\begin{minted}{python}
"""OpenAPI breaking-change gate (Chapter 9, Listing 9.2).

Classifies the delta between an OLD and a NEW OpenAPI 3.x YAML spec as
BREAKING or NON-BREAKING and exits non-zero when any breaking change exists.
Stdlib + PyYAML only. Inline schemas assumed (no $ref resolution).

Usage (PowerShell):
    python openapi_diff_gate.py catalog_v1.yaml catalog_v2.yaml
"""

from __future__ import annotations

import argparse
import enum
import sys
from dataclasses import dataclass
from pathlib import Path
from typing import Any

import yaml

HTTP_METHODS = frozenset(
    {"get", "put", "post", "delete", "patch", "head", "options", "trace"}
)


class Severity(enum.Enum):
    BREAKING = "BREAKING"
    NON_BREAKING = "NON-BREAKING"


@dataclass(frozen=True)
class Change:
    severity: Severity
    location: str
    message: str

    def render(self) -> str:
        return f"[{self.severity.value:>12}] {self.location}: {self.message}"


def _operations(path_item: dict[str, Any]) -> dict[str, Any]:
    return {m: op for m, op in path_item.items() if m.lower() in HTTP_METHODS}


def _diff_schema(
    old: dict[str, Any],
    new: dict[str, Any],
    location: str,
    direction: str,  # "request" or "response"
    changes: list[Change],
) -> None:
    old_props: dict[str, Any] = old.get("properties", {}) or {}
    new_props: dict[str, Any] = new.get("properties", {}) or {}
    new_required = set(new.get("required", []) or [])

    for name in old_props:
        if name not in new_props:
            changes.append(Change(
                Severity.BREAKING, location,
                f"property '{name}' removed from {direction} schema"))

    for name, new_prop in new_props.items():
        if name not in old_props:
            if direction == "request" and name in new_required:
                changes.append(Change(
                    Severity.BREAKING, location,
                    f"required request property '{name}' added"))
            else:
                changes.append(Change(
                    Severity.NON_BREAKING, location,
                    f"optional {direction} property '{name}' added"))
            continue
        old_type = old_props[name].get("type")
        new_type = new_prop.get("type")
        if old_type != new_type:
            changes.append(Change(
                Severity.BREAKING, location,
                f"property '{name}' type changed {old_type!r} -> {new_type!r}"))
        elif old_type == "object":
            _diff_schema(old_props[name], new_prop,
                         f"{location}.{name}", direction, changes)
        elif old_type == "array":
            _diff_schema(old_props[name].get("items", {}) or {},
                         new_prop.get("items", {}) or {},
                         f"{location}.{name}[]", direction, changes)
        old_enum = old_props[name].get("enum")
        new_enum = new_prop.get("enum")
        if old_enum and new_enum:
            removed = set(old_enum) - set(new_enum)
            if removed:
                changes.append(Change(
                    Severity.BREAKING, location,
                    f"enum value(s) removed from '{name}': {sorted(removed)}"))
            added = set(new_enum) - set(old_enum)
            if added:
                changes.append(Change(
                    Severity.NON_BREAKING, location,
                    f"enum value(s) added to '{name}': {sorted(added)} "
                    "(safe only if the enum is documented as open)"))


def _json_schema(media: dict[str, Any]) -> dict[str, Any]:
    content = media.get("content", {}) or {}
    return (content.get("application/json", {}) or {}).get("schema", {}) or {}


def _diff_operation(
    old: dict[str, Any],
    new: dict[str, Any],
    location: str,
    changes: list[Change],
) -> None:
    old_body = _json_schema(old.get("requestBody", {}) or {})
    new_body = _json_schema(new.get("requestBody", {}) or {})
    if new_body and old_body:
        _diff_schema(old_body, new_body, f"{location} requestBody",
                     "request", changes)
    elif new_body and not old_body and new_body.get("required"):
        changes.append(Change(
            Severity.BREAKING, location,
            "request body with required properties introduced"))

    old_responses: dict[str, Any] = old.get("responses", {}) or {}
    new_responses: dict[str, Any] = new.get("responses", {}) or {}
    for code in old_responses:
        if code not in new_responses:
            changes.append(Change(
                Severity.BREAKING, location,
                f"response status '{code}' removed"))
    for code in new_responses:
        if code not in old_responses:
            changes.append(Change(
                Severity.NON_BREAKING, location,
                f"response status '{code}' added"))
    for code in old_responses.keys() & new_responses.keys():
        _diff_schema(_json_schema(old_responses[code]),
                     _json_schema(new_responses[code]),
                     f"{location} response {code}", "response", changes)


def diff_specs(old: dict[str, Any], new: dict[str, Any]) -> list[Change]:
    changes: list[Change] = []
    old_paths: dict[str, Any] = old.get("paths", {}) or {}
    new_paths: dict[str, Any] = new.get("paths", {}) or {}

    for path in old_paths:
        if path not in new_paths:
            changes.append(Change(Severity.BREAKING, path, "path removed"))
    for path in new_paths:
        if path not in old_paths:
            changes.append(Change(Severity.NON_BREAKING, path, "path added"))

    for path in old_paths.keys() & new_paths.keys():
        old_ops = _operations(old_paths[path])
        new_ops = _operations(new_paths[path])
        for method in old_ops:
            if method not in new_ops:
                changes.append(Change(
                    Severity.BREAKING, f"{method.upper()} {path}",
                    "operation removed"))
        for method in new_ops:
            if method not in old_ops:
                changes.append(Change(
                    Severity.NON_BREAKING, f"{method.upper()} {path}",
                    "operation added"))
        for method in old_ops.keys() & new_ops.keys():
            _diff_operation(old_ops[method], new_ops[method],
                            f"{method.upper()} {path}", changes)
    return changes


def main(argv: list[str] | None = None) -> int:
    parser = argparse.ArgumentParser(description=__doc__)
    parser.add_argument("old_spec", type=Path)
    parser.add_argument("new_spec", type=Path)
    args = parser.parse_args(argv)

    old = yaml.safe_load(args.old_spec.read_text(encoding="utf-8"))
    new = yaml.safe_load(args.new_spec.read_text(encoding="utf-8"))
    changes = diff_specs(old, new)

    breaking = [c for c in changes if c.severity is Severity.BREAKING]
    for change in changes:
        print(change.render())
    print(f"\n{len(changes)} change(s): {len(breaking)} breaking, "
          f"{len(changes) - len(breaking)} non-breaking")
    if breaking:
        print("GATE RESULT: FAIL -- breaking changes require a new version")
        return 1
    print("GATE RESULT: PASS -- all changes are backward compatible")
    return 0


if __name__ == "__main__":
    sys.exit(main())
\end{minted}

The two fixture specs below describe the Catalog contract before and after
the breaking v1-to-v2 redesign. Save them beside the gate.

\begin{minted}{yaml}
# catalog_v1.yaml -- published contract (the baseline the gate protects)
openapi: "3.1.0"
info:
  title: Northwind Robotics Catalog API
  version: "1.4.0"
paths:
  /catalog:
    get:
      operationId: listCatalog
      responses:
        "200":
          description: All parts
          content:
            application/json:
              schema:
                type: object
                required: [items]
                properties:
                  items:
                    type: array
                    items:
                      type: object
                      required: [id, sku, name, price, category, in_stock]
                      properties:
                        id: {type: string}
                        sku: {type: string}
                        name: {type: string}
                        price: {type: number}
                        category: {type: string, enum: [actuator, sensor, controller]}
                        in_stock: {type: boolean}
  /catalog/{partId}:
    get:
      operationId: getPart
      responses:
        "200":
          description: One part
          content:
            application/json:
              schema:
                type: object
                required: [id, sku, name, price, category, in_stock]
                properties:
                  id: {type: string}
                  sku: {type: string}
                  name: {type: string}
                  price: {type: number}
                  category: {type: string, enum: [actuator, sensor, controller]}
                  in_stock: {type: boolean}
        "404":
          description: Not found
\end{minted}

\begin{minted}{yaml}
# catalog_v2.yaml -- proposed contract (renames and retypes => gate must FAIL)
openapi: "3.1.0"
info:
  title: Northwind Robotics Catalog API
  version: "2.0.0"
paths:
  /catalog:
    get:
      operationId: listCatalog
      responses:
        "200":
          description: All parts
          content:
            application/json:
              schema:
                type: object
                required: [items]
                properties:
                  items:
                    type: array
                    items:
                      type: object
                      required: [part_id, sku, name, unit_price_cents, category, stock]
                      properties:
                        part_id: {type: string}
                        sku: {type: string}
                        name: {type: string}
                        unit_price_cents: {type: integer}
                        currency: {type: string}
                        category: {type: string, enum: [actuator, sensor, controller, battery]}
                        stock: {type: integer}
  /catalog/{partId}:
    get:
      operationId: getPart
      responses:
        "200":
          description: One part
          content:
            application/json:
              schema:
                type: object
                required: [part_id, sku, name, unit_price_cents, category, stock]
                properties:
                  part_id: {type: string}
                  sku: {type: string}
                  name: {type: string}
                  unit_price_cents: {type: integer}
                  currency: {type: string}
                  category: {type: string, enum: [actuator, sensor, controller, battery]}
                  stock: {type: integer}
        "404":
          description: Not found
  /catalog/health:
    get:
      operationId: catalogHealth
      responses:
        "200":
          description: Liveness probe
\end{minted}

\subsection{Listing 9.3 --- Media-Type Version Negotiator}
\label{subsec:ch9-listing3}

A correct negotiator for
\texttt{Accept: application/vnd.northwind.catalog+json; version=2}: it parses
q-values and parameters per RFC~9110~\cite{rfc9110}, selects the best
supported version, and yields 406 (with a \texttt{Link} to supported
versions) when nothing matches.

\begin{minted}{python}
"""Media-type version negotiation (Chapter 9, Listing 9.3).

Parses Accept headers carrying vendor media types with a `version=N`
parameter, honors q-values, and returns None (-> HTTP 406) when no supported
version matches. Run the module for a deterministic negotiation table plus a
TestClient demonstration. Stdlib + FastAPI.
"""

from __future__ import annotations

from dataclasses import dataclass

from fastapi import FastAPI, Request
from fastapi.responses import JSONResponse
from fastapi.testclient import TestClient

VENDOR_TYPE = "application/vnd.northwind.catalog+json"
SUPPORTED_VERSIONS = (1, 2)
DEFAULT_VERSION = 2

PART_V1 = {"id": "WR-1042", "sku": "SRV-T450", "name": "Torque Servo T450",
           "price": 129.99, "category": "actuator", "in_stock": True}
PART_V2 = {"part_id": "WR-1042", "sku": "SRV-T450", "name": "Torque Servo T450",
           "unit_price_cents": 12999, "currency": "USD",
           "category": "actuator", "stock": 34}


@dataclass(frozen=True)
class MediaRange:
    media_type: str
    q: float
    version: int | None
    order: int


def parse_accept(header: str | None) -> list[MediaRange]:
    """Parse an Accept header into q-sorted media ranges (stable on ties)."""

    if not header:
        return [MediaRange("*/*", 1.0, None, 0)]
    ranges: list[MediaRange] = []
    for order, raw in enumerate(header.split(",")):
        tokens = [t.strip() for t in raw.split(";")]
        media_type = tokens[0].lower()
        q, version = 1.0, None
        for param in tokens[1:]:
            if "=" not in param:
                continue
            key, _, value = param.partition("=")
            key, value = key.strip().lower(), value.strip().strip('"')
            if key == "q":
                try:
                    q = float(value)
                except ValueError:
                    q = 0.0
            elif key == "version":
                try:
                    version = int(value)
                except ValueError:
                    version = None
        ranges.append(MediaRange(media_type, q, version, order))
    acceptable = [r for r in ranges if r.q > 0.0]
    return sorted(acceptable, key=lambda r: (-r.q, r.order))


def negotiate(
    header: str | None,
    supported: tuple[int, ...] = SUPPORTED_VERSIONS,
    default: int = DEFAULT_VERSION,
) -> int | None:
    """Return the best supported version for the header, or None (-> 406).

    A vendor range with an explicit version matches only that version; a bare
    vendor type, application/json, or */* maps to the server default.
    """

    best_version: int | None = None
    best_q = -1.0
    for r in parse_accept(header):
        if r.media_type == VENDOR_TYPE and r.version is not None:
            if r.version in supported and r.q > best_q:
                best_version, best_q = r.version, r.q
        elif r.media_type in (VENDOR_TYPE, "application/json", "*/*"):
            if r.q > best_q:
                best_version, best_q = default, r.q
    return best_version


app = FastAPI(title="Northwind Robotics Catalog (media-type versioned)")


@app.get("/catalog/{part_id}")
def get_part(part_id: str, request: Request) -> JSONResponse:
    version = negotiate(request.headers.get("accept"))
    if version is None:
        return JSONResponse(
            status_code=406,
            content={"type": "about:blank", "title": "Not Acceptable",
                     "status": 406,
                     "detail": "supported: " + ", ".join(
                         f"{VENDOR_TYPE}; version={v}"
                         for v in SUPPORTED_VERSIONS)},
            headers={"Link": f'</catalog>; rel="versions"'},
        )
    body = PART_V1 if version == 1 else PART_V2
    return JSONResponse(
        content=body,
        media_type=f"{VENDOR_TYPE}; version={version}",
        headers={"Vary": "Accept"},
    )


def demo() -> None:
    cases = [
        None,
        "application/json",
        "application/vnd.northwind.catalog+json; version=1",
        "application/vnd.northwind.catalog+json; version=2",
        "application/vnd.northwind.catalog+json; version=1;q=0.4, "
        "application/vnd.northwind.catalog+json; version=2;q=0.9",
        "application/vnd.northwind.catalog+json; version=3",
        "application/xml",
    ]
    print(f"{'Accept header':<78} -> version")
    for header in cases:
        print(f"{str(header):<78} -> {negotiate(header)}")
    client = TestClient(app)
    r = client.get("/catalog/WR-1042",
                   headers={"Accept": f"{VENDOR_TYPE}; version=3"})
    print("\n406 case:", r.status_code, r.json()["detail"])


if __name__ == "__main__":
    demo()
\end{minted}

\subsection{Listing 9.4 --- Consumer-Tolerance Test Suite}
\label{subsec:ch9-listing4}

Compatibility is a property you \emph{test}, not assert. This pytest suite
plays a frozen v1-era client against the live dual-serving app of
Listing~9.1 and proves the tolerant-reader guarantees: unknown fields are
ignored, additive changes are invisible, new enum values degrade to a
fallback instead of crashing, and the v1 surface keeps its promised shape
plus its deprecation metadata.

\begin{minted}{python}
"""Consumer-tolerance contract tests (Chapter 9, Listing 9.4).

Proves that a frozen v1-era client keeps working while the API evolves:
unknown-field tolerance, additive changes, open-enum fallback, and the
stability of the deprecated v1 surface. Run: pytest -q
"""

from __future__ import annotations

import pytest
from fastapi.testclient import TestClient

from catalog_api import app

# A v1-era client knows exactly these fields and nothing else.
LEGACY_FIELDS = frozenset({"id", "sku", "name", "price", "category", "in_stock"})
LEGACY_CATEGORIES = frozenset({"actuator", "sensor", "controller"})


def legacy_parse_part(payload: dict) -> dict:
    """Tolerant reader: extract known fields, ignore everything else."""

    return {k: payload[k] for k in LEGACY_FIELDS & payload.keys()}


def legacy_category(value: str) -> str:
    """Open-enum fallback rule documented in the v1 contract."""

    return value if value in LEGACY_CATEGORIES else "unknown"


@pytest.fixture(scope="module")
def client() -> TestClient:
    return TestClient(app)


def test_v1_surface_keeps_every_promised_field(client: TestClient) -> None:
    response = client.get("/v1/catalog")
    assert response.status_code == 200
    parts = response.json()
    assert parts, "fixture catalog must not be empty"
    for part in parts:
        assert LEGACY_FIELDS <= set(part), f"v1 contract violated: {part}"
        assert isinstance(part["price"], float)
        assert isinstance(part["in_stock"], bool)


def test_unknown_fields_are_ignored_by_tolerant_reader() -> None:
    future_payload = {
        "id": "WR-9000", "sku": "NXT-1", "name": "Next-Gen Servo",
        "price": 42.5, "category": "actuator", "in_stock": True,
        # Fields a future version might add:
        "unit_price_cents": 4250, "stock": 7, "tags": ["servo"],
        "dimensions": {"w_mm": 40, "h_mm": 20}, "carbon_footprint_g": 812,
    }
    parsed = legacy_parse_part(future_payload)
    assert parsed == {
        "id": "WR-9000", "sku": "NXT-1", "name": "Next-Gen Servo",
        "price": 42.5, "category": "actuator", "in_stock": True,
    }


def test_new_enum_value_degrades_to_fallback_not_crash() -> None:
    # 'battery' was added after this client shipped; fallback keeps it alive.
    assert legacy_category("battery") == "unknown"
    assert legacy_category("actuator") == "actuator"


def test_v2_additive_fields_do_not_leak_into_v1(client: TestClient) -> None:
    v1 = client.get("/v1/catalog/WR-1042").json()
    v2 = client.get("/v2/catalog/WR-1042").json()
    assert set(v1) == LEGACY_FIELDS
    assert {"unit_price_cents", "stock", "currency", "tags"} <= set(v2)
    # Adapter fidelity: v1 values are derivable from the v2 model.
    assert v1["price"] == round(v2["unit_price_cents"] / 100, 2)
    assert v1["in_stock"] == (v2["stock"] > 0)


def test_v1_carries_deprecation_metadata_v2_does_not(client: TestClient) -> None:
    v1 = client.get("/v1/catalog/WR-1042")
    assert v1.headers["Deprecation"].startswith("@")
    assert "Sunset" in v1.headers
    assert 'rel="successor-version"' in v1.headers["Link"]
    v2 = client.get("/v2/catalog/WR-1042")
    assert "Deprecation" not in v2.headers
    assert "Sunset" not in v2.headers
\end{minted}

\subsection{Listing 9.5 --- Retirement-Readiness Report Generator}
\label{subsec:ch9-listing5}

Retirement is a data-driven decision. This tool consumes a JSON-lines
access-log fixture (gateway logs exported offline), computes per-version
traffic share, distinct consumers, and monthly trend, and emits a verdict
against a threshold policy. The embedded fixture shows v1 declining toward
the 5\%\ retirement line.

\begin{minted}{python}
"""Retirement-readiness report generator (Chapter 9, Listing 9.5).

Reads a JSON-lines access log (one event per line: ts, path, consumer) and
reports per-version usage share, distinct consumers, and monthly trend, then
emits a retirement verdict against the chapter's threshold policy.

Usage (PowerShell):
    python retirement_report.py --demo            # embedded fixture
    python retirement_report.py access_log.jsonl  # your own export
"""

from __future__ import annotations

import argparse
import json
from collections import defaultdict
from pathlib import Path

RETIRE_SHARE_THRESHOLD = 0.05  # retire when v1 < 5% of catalog traffic...

DEMO_LOG = """\
{"ts": "2026-05-04T09:13:00Z", "path": "/v1/catalog/WR-1042", "consumer": "partner-acme"}
{"ts": "2026-05-05T10:01:00Z", "path": "/v2/catalog", "consumer": "storefront-web"}
{"ts": "2026-05-06T11:40:00Z", "path": "/v2/catalog/WR-2210", "consumer": "storefront-web"}
{"ts": "2026-06-02T09:00:00Z", "path": "/v2/catalog", "consumer": "storefront-web"}
{"ts": "2026-06-03T09:05:00Z", "path": "/v2/catalog", "consumer": "partner-globex"}
{"ts": "2026-06-04T12:31:00Z", "path": "/v2/catalog/WR-4100", "consumer": "warehouse-wms"}
{"ts": "2026-06-05T15:10:00Z", "path": "/v2/catalog", "consumer": "storefront-web"}
{"ts": "2026-07-01T09:00:00Z", "path": "/v2/catalog", "consumer": "storefront-web"}
{"ts": "2026-07-02T09:10:00Z", "path": "/v2/catalog/WR-3307", "consumer": "partner-globex"}
{"ts": "2026-07-03T10:20:00Z", "path": "/v2/catalog", "consumer": "warehouse-wms"}
{"ts": "2026-07-04T11:00:00Z", "path": "/v2/catalog", "consumer": "mobile-ios-4.0"}
{"ts": "2026-07-05T13:45:00Z", "path": "/v2/catalog/WR-1042", "consumer": "storefront-web"}
{"ts": "2026-08-02T08:30:00Z", "path": "/v2/catalog", "consumer": "storefront-web"}
{"ts": "2026-08-02T09:00:00Z", "path": "/v2/catalog", "consumer": "partner-globex"}
{"ts": "2026-08-03T09:30:00Z", "path": "/v2/catalog/WR-4100", "consumer": "warehouse-wms"}
{"ts": "2026-08-04T10:00:00Z", "path": "/v2/catalog", "consumer": "mobile-ios-4.0"}
{"ts": "2026-08-05T11:15:00Z", "path": "/v2/catalog/WR-2210", "consumer": "storefront-web"}
{"ts": "2026-08-06T12:00:00Z", "path": "/v2/catalog", "consumer": "partner-globex"}
{"ts": "2026-08-07T12:30:00Z", "path": "/v2/catalog/WR-3307", "consumer": "warehouse-wms"}
{"ts": "2026-08-08T13:00:00Z", "path": "/v2/catalog", "consumer": "mobile-ios-4.0"}
{"ts": "2026-08-09T14:20:00Z", "path": "/v2/catalog/WR-1042", "consumer": "storefront-web"}
"""


def path_version(path: str) -> str | None:
    for prefix in ("/v1/", "/v2/"):
        if path.startswith(prefix):
            return prefix.strip("/")
    return None


def build_report(lines: list[str]) -> str:
    total: dict[str, int] = defaultdict(int)
    consumers: dict[str, set[str]] = defaultdict(set)
    monthly: dict[str, dict[str, int]] = defaultdict(lambda: defaultdict(int))

    for raw in lines:
        raw = raw.strip()
        if not raw:
            continue
        event = json.loads(raw)
        version = path_version(str(event["path"]))
        if version is None:
            continue  # unversioned traffic is out of scope for this report
        total[version] += 1
        consumers[version].add(str(event["consumer"]))
        monthly[str(event["ts"])[:7]][version] += 1

    grand_total = sum(total.values())
    if grand_total == 0:
        return "No versioned traffic found; nothing to report."

    out = ["RETIREMENT-READINESS REPORT -- Northwind Robotics Catalog API",
           "=" * 62, f"Total versioned requests: {grand_total}", ""]
    for version in sorted(total):
        share = total[version] / grand_total
        out.append(f"{version}: {total[version]:>4} requests "
                   f"({share:6.1%}), {len(consumers[version])} consumer(s): "
                   f"{sorted(consumers[version])}")
    out += ["", "Monthly trend (v1 share): "]
    for month in sorted(monthly):
        m = monthly[month]
        month_total = sum(m.values())
        out.append(f"  {month}: v1={m.get('v1', 0):>2} v2={m.get('v2', 0):>2} "
                   f"(v1 share {m.get('v1', 0) / month_total:6.1%})")

    months = sorted(monthly)
    first = monthly[months[0]]
    last = monthly[months[-1]]
    first_share = first.get("v1", 0) / sum(first.values())
    last_share = last.get("v1", 0) / sum(last.values())
    declining = last_share <= first_share
    v1_share = total.get("v1", 0) / grand_total

    out += ["", "Verdict:"]
    if v1_share < RETIRE_SHARE_THRESHOLD and declining:
        out.append(f"  v1 share {v1_share:.1%} < {RETIRE_SHARE_THRESHOLD:.0%} "
                   "and declining: SCHEDULE SUNSET BROWNOUTS, confirm the "
                   f"remaining consumer(s) {sorted(consumers.get('v1', []))}, "
                   "then retire on the published date.")
    else:
        out.append(f"  v1 share {v1_share:.1%} (declining={declining}): KEEP "
                   "v1 RUNNING; continue migration campaign and re-evaluate "
                   "next month.")
    return "\n".join(out)


def main(argv: list[str] | None = None) -> int:
    parser = argparse.ArgumentParser(description=__doc__)
    parser.add_argument("logfile", type=Path, nargs="?",
                        help="JSON-lines access log export")
    parser.add_argument("--demo", action="store_true",
                        help="use the embedded deterministic fixture")
    args = parser.parse_args(argv)

    if args.demo or args.logfile is None:
        lines = DEMO_LOG.splitlines()
    else:
        lines = args.logfile.read_text(encoding="utf-8").splitlines()
    print(build_report(lines))
    return 0


if __name__ == "__main__":
    raise SystemExit(main())
\end{minted}

%% ----------------------------------------------------------------------------
\section{Common Pitfalls \& Anti-Patterns}
\label{sec:ch9-pitfalls}

\begin{warningbox}
Every anti-pattern below has caused a real production incident class. They
are ordered roughly by frequency of appearance in design reviews.
\end{warningbox}

\paragraph{1. Version-per-endpoint chaos.}
Letting individual endpoints evolve their own versions independently
(\texttt{/v1/orders}, \texttt{/v3/customers}, \texttt{/v2/orders/items})
shatters the API into a combinatorial mess: no single document describes the
system, client SDKs must track a version matrix, and gateway routing becomes
archaeology. Version the \emph{API} (or at most the bounded context), not the
endpoint. Media-type versioning is the one legitimate exception because its
per-resource granularity is negotiated, not routed.

\paragraph{2. Silently changing semantics within a version.}
Renaming nothing, retyping nothing --- just quietly redefining what
\texttt{stock} means from ``on hand'' to ``available to promise.'' No client
errors fire; every client silently computes wrong answers. Semantic drift is
the most damaging break precisely because nothing detects it except business
outcomes. The defense is contractual: field \emph{meaning} is part of the
contract, documented, and covered by consumer-tolerance tests with fixed
fixtures.

\paragraph{3. Breaking changes via ``harmless'' renames.}
\texttt{price} becomes \texttt{unit\_price\_cents}; \texttt{id} becomes
\texttt{part\_id}. Each rename is internally justified (clarity! units in the
name!) and externally a removal-plus-addition --- a breaking change wearing a
refactor's clothes (Listing~9.2's fixture pair demonstrates exactly this, and
the gate fails the build). If a name must change, add the new name, keep
serving the old one, and migrate.

\paragraph{4. Never retiring versions.}
Versions without sunsets accumulate into permanent carrying cost
(Equation~\ref{eq:ch9-cost}): every adapter, doc set, and CI lane lives
forever, and eventually the team maintains contracts nobody remembers
designing. A version without a deprecation plan is a mortgage with no
amortization schedule.

\paragraph{5. Version in both URI and header.}
\texttt{/v1/catalog} with \texttt{X-API-Version: 2} --- which wins? Every
consumer picks a different answer, support tickets multiply, and the gateway
must implement a precedence rule that no document defined. Choose exactly one
version signal per API and reject requests that carry conflicting ones.

\paragraph{6. Ignoring cache poisoning across versions.}
With header- or media-type-based versioning, a shared cache keyed on URI
alone will happily serve a v1 body to a v2 request. The mitigations are
mandatory, not optional: emit \texttt{Vary: X-API-Version} (or
\texttt{Vary: Accept}) on every versioned response and verify intermediaries
honor it~\cite{rfc9110} --- or use URI versioning and let the URI do the
keying. Listing~9.3 emits \texttt{Vary: Accept} for this reason.

%% ----------------------------------------------------------------------------
\section{Hands-On Production Lab \& Real-World Case Study}
\label{sec:ch9-lab}

\subsection*{Lab: Evolve the Catalog API from v1 to v2}
You own the Northwind Robotics Catalog API. The v1 contract (float-dollar
\texttt{price}, boolean \texttt{in\_stock}) cannot express multi-currency
pricing or real stock counts, so a v2 contract has been designed. Execute the
full expand-and-contract lifecycle.

\textbf{Setup (PowerShell):}
\begin{minted}{text}
cd C:\Training\ch09
python -m venv .venv
.\.venv\Scripts\Activate.ps1
pip install -r requirements.txt   # see Section 9.11
\end{minted}

\begin{enumerate}
  \item \textbf{Expand.} Save Listing~9.1 as \texttt{catalog\_api.py} and run
        \texttt{python catalog\_api.py}. Observe that v1 and v2 both answer
        and that v1 responses carry \texttt{Deprecation}, \texttt{Sunset},
        and \texttt{Link} headers while v2 responses do not.
  \item \textbf{Gate the change.} Save Listing~9.2 as
        \texttt{openapi\_diff\_gate.py} plus the two fixture specs and run
        \texttt{python openapi\_diff\_gate.py catalog\_v1.yaml catalog\_v2.yaml};
        confirm it reports the renames and retypes as BREAKING and exits 1
        (\texttt{echo \$LASTEXITCODE}). Then diff \texttt{catalog\_v1.yaml}
        against itself and confirm the gate passes and exits 0 --- this is
        the ``no changes'' CI lane.
  \item \textbf{Prove consumer survival.} Save Listing~9.4 as
        \texttt{test\_consumer\_tolerance.py} and run
        \texttt{pytest -q}. All five tests must pass: they are the evidence
        that the expand phase harmed no one.
  \item \textbf{Negotiate.} Save Listing~9.3 as
        \texttt{media\_type\_negotiation.py}, run it, and verify the
        negotiation table: explicit versions select their representation,
        bare types get the default, unsupported versions yield \texttt{None}
        (HTTP 406).
  \item \textbf{Contract.} Save Listing~9.5 as
        \texttt{retirement\_report.py} and run
        \texttt{python retirement\_report.py --demo}. Read the verdict: v1 is
        below the 5\%\ threshold and declining, with exactly one identified
        straggler --- the campaign's final contact list.
\end{enumerate}

\subsection*{War Story: The ``Harmless'' Enum Addition}
A robotics parts platform (the kind Northwind Robotics sells into) shipped a
seemingly trivial release: the catalog \texttt{category} enum gained a new
value, \texttt{battery}, to support a product-line launch. The API team ran
their compatibility checks --- no fields removed, no types changed, all
endpoints identical --- and deployed on a Tuesday. By Wednesday morning,
crash dashboards for the two largest mobile consumers were vertical. Both
apps had been built from generated models in which the enum was a
\emph{closed} native type; their JSON decoders threw on the unknown value,
and because decoding happened inside the list parser, a single battery part
anywhere in a page of results poisoned the entire response. The apps did not
degrade --- they died on launch screens that happened to render category
chips.

The postmortem lessons are this chapter in miniature. (1)~The producer's
compatibility model (additive = safe) did not match the consumers' parsing
model (closed enum = fatal), and under Definition~\ref{def:breaking} the
change \emph{was} breaking for that consumer base. (2)~The enum had never
been documented as open, so consumers reasonably treated it as closed.
(3)~The fix was not code but contract: the API now publishes every enum as
extensible with an \texttt{unknown} sentinel, consumer SDKs regenerate with
fallback arms, a consumer-tolerance fixture suite (the pattern of
Listing~9.4) injects synthetic unknown enum values in CI, and genuinely
closed enums can only change in new major versions. The enum addition that
crashed the fleet is now the first slide of their API design onboarding.

%% ----------------------------------------------------------------------------
\section{Self-Assessment Exercises \& Deep-Dive Questions}
\label{sec:ch9-exercises}

\begin{enumerate}[label=\textbf{E9.\arabic*.}, leftmargin=2.2cm]
  \item \textbf{Classification drill.} Classify each as breaking or
        non-breaking, citing Definition~\ref{def:breaking}: (a) adding an
        optional \texttt{warehouse} field to the part response; (b) making a
        previously optional request field required; (c) changing
        \texttt{price} from dollars to cents while renaming it
        \texttt{unit\_price\_cents}; (d) reordering JSON object members;
        (e) changing a 404 detail message; (f) adding \texttt{410 Gone} as a
        \emph{new} response for retired resources.
  \item \textbf{Extend the gate.} Listing~9.2 ignores path and query
        \texttt{parameters}. Extend it so that adding a \emph{required}
        parameter or removing an existing parameter is flagged breaking.
        Add fixtures proving both directions.
  \item \textbf{q-value correctness.} Construct an \texttt{Accept} header in
        which a lower-q explicit version must lose to a higher-q wildcard,
        and verify Listing~9.3's \texttt{negotiate} returns the default.
        Then construct the reverse.
  \item \textbf{Adapter audit.} Propose a v2 field that \emph{cannot} be
        back-translated into the v1 surface of Listing~9.1 (e.g., a
        multi-currency price list). What are the three honest options, and
        which does expand-and-contract prescribe?
  \item \textbf{Cache audit.} Your API uses \texttt{X-API-Version}. Write the
        exact response headers required for safe shared caching, and explain
        what breaks if a CDN ignores \texttt{Vary}.
  \item \textbf{Deep dive: closed ecosystems.} Your consumers are all mobile
        apps with six-month release laggards and generated strict models.
        Defend or reject media-type versioning for this base, and choose
        deprecation timelines accordingly.
  \item \textbf{Deep dive: the exception clause.} Your threat model team
        demands an immediate breaking change to close an auth bypass.
        Reconcile ``never break within a version'' with security response;
        define the policy you would publish \emph{before} the incident.
\end{enumerate}

%% ----------------------------------------------------------------------------
\section{Interview Q\&A Vault}
\label{sec:ch9-interview}

\subsection*{Junior (Q1--Q10)}
\begin{enumerate}[label=\textbf{Q\arabic*.}, leftmargin=1.6cm]
  \item \textbf{What is an API version?}
        A distinctly identified, supported contract between the API and its
        consumers --- not just a URL prefix. Minting one declares that the
        contract changed incompatibly; supporting one commits you to its
        stability SLA.
  \item \textbf{Define a breaking change.}
        Any change after which at least one formerly valid, successful client
        interaction is rejected, fails, or becomes uninterpretable by a
        client written against the old contract (Definition~\ref{def:breaking}).
  \item \textbf{Classify these five changes: removing a response field;
        adding an optional request field; renaming \texttt{price} to
        \texttt{unit\_price}; adding a new endpoint; changing a field from
        float to integer.}
        Breaking; non-breaking; breaking (rename = remove + add);
        non-breaking; breaking.
  \item \textbf{Is adding a response field breaking?}
        For tolerant readers, no --- unknown members are ignored. For clients
        with strict schema validation or closed generated models, it can be.
        Your contract must state which client posture you support.
  \item \textbf{Why is URI versioning so popular despite being ``un-RESTful''?}
        Operational clarity: visible in logs and browsers, trivially routed
        by gateways, naturally cache-keyed, and effortless for consumers to
        pin. Most organizations value those over URI purity.
  \item \textbf{What does HTTP 406 mean in a versioned API?}
        Not Acceptable --- the server cannot produce any representation the
        client's \texttt{Accept} header permits, e.g.\ an unsupported
        \texttt{version=N} media-type parameter.
  \item \textbf{What is the difference between deprecation and retirement?}
        Deprecated means still functional but discouraged, no longer
        enhanced, and scheduled for removal; retired (sunset) means the
        surface no longer responds (typically 410 Gone or a redirect).
  \item \textbf{What are the \texttt{Deprecation} and \texttt{Sunset}
        headers?}
        Standardized response headers (RFC~8594, with \texttt{Deprecation}
        revised by RFC~9745) announcing when a surface was deprecated and
        when it will become unresponsive, so clients can detect lifecycle
        state programmatically~\cite{rfc8594}.
  \item \textbf{Give two examples of non-breaking changes.}
        Adding an endpoint; adding an optional request field or a response
        field (for tolerant readers); adding an enum value to a
        documented-open enum; loosening validation.
  \item \textbf{Why should you never reuse a retired field name?}
        Stored payloads, caches, and old clients still carry the old meaning;
        reusing the name silently reinterprets their data. Retired names are
        reserved forever --- Protobuf enforces this with \texttt{reserved}
        \cite{protobufdocs}.
\end{enumerate}

\subsection*{Mid (Q11--Q20)}
\begin{enumerate}[label=\textbf{Q\arabic*.}, leftmargin=1.6cm, start=11]
  \item \textbf{URI vs.\ header versioning --- defend a choice for a public
        partner API.}
        URI path: partners debug with curl and browsers, gateways route on
        paths for free, caches key correctly without \texttt{Vary}, and
        pinning is obvious in code review. I would accept the REST-purity
        cost and document why. (Full credit answers argue either side with
        the trade-off matrix, not dogma.)
  \item \textbf{How does Stripe's date-based versioning work?}
        Accounts pin an API version by date; requests may override per call
        with a version header. Internally, responses are rendered from
        current code through a stack of backward-compatibility transforms
        down to the pinned date, so integrations never change underneath a
        consumer until they choose to upgrade.
  \item \textbf{What is expand-and-contract (parallel change)?}
        A three-phase migration: \emph{expand} (add the new contract beside
        the old, remove nothing), \emph{migrate} (move consumers with
        telemetry and campaigns), \emph{contract} (retire the old surface
        only when usage proves it safe). It keeps an intermediate state where
        both contracts are valid, so consumers never see a discontinuity
        \cite{kleppmann2017ddia}.
  \item \textbf{What is an anti-corruption layer in versioned APIs?}
        An adapter that translates between an old public contract and the
        current domain model, preventing legacy vocabulary from leaking into
        new code. It must be pure, total, and contract-tested; fields that
        cannot be faithfully derived signal retirement, not more translation.
  \item \textbf{How do you dual-serve two versions without duplicating
        business logic?}
        One service core owns the newest domain model; older versions are
        thin adapters over it; a gateway routes by version signal; telemetry
        on the router feeds the retirement decision
        (Figure~\ref{fig:ch9-dual}).
  \item \textbf{How do caches interact with header-based versioning?}
        Caches key primarily on URI; two versions sharing a URI poison each
        other unless responses carry \texttt{Vary} on the versioning header
        and every intermediary honors it~\cite{rfc9110}. URI versioning
        sidesteps the problem entirely.
  \item \textbf{When may you break compatibility without minting a new
        version?}
        Only under pre-published exception clauses: critical security fixes,
        legally mandated changes, correcting behavior that contradicts the
        documented contract, or surfaces explicitly labeled unstable/beta.
        The clause must exist \emph{before} it is invoked.
  \item \textbf{Your consumer base can't tolerate new enum values. What are
        your options?}
        Treat enum growth as breaking: mint a new version; or add a parallel
        string field with the extended value set and freeze the enum; or run
        a migration campaign to tolerant SDKs first, then grow the enum.
        What you may not do is add the value and hope.
  \item \textbf{What belongs in a migration guide?}
        A field-level old-to-new mapping table, behavioral deltas (semantics,
        errors, defaults), copy-pasteable request/response diffs, SDK upgrade
        steps, the sunset timeline, and a support channel.
  \item \textbf{How does SemVer apply to APIs?}
        MAJOR = breaking contract change (new version); MINOR =
        backward-compatible addition (field, endpoint, open-enum value);
        PATCH = compatible fix. The changelog classifies every dated entry
        by this grammar.
\end{enumerate}

\subsection*{Senior (Q21--Q27)}
\begin{enumerate}[label=\textbf{Q\arabic*.}, leftmargin=1.6cm, start=21]
  \item \textbf{Design a breaking-change CI gate.}
        Diff the proposed OpenAPI description against the published baseline;
        classify deltas (removed paths/operations/properties, type changes,
        new required request fields, removed enum values) as breaking; fail
        the pipeline unless the PR also bumps the major version. Keep the
        classifier inspectable (Listing~9.2) and treat warnings (enum
        additions) as human-review triggers~\cite{openapi31}.
  \item \textbf{How do you decide when to retire a version?}
        Telemetry, not calendar: schedule sunset when the old version falls
        below a traffic-share threshold (commonly $\sim$5\%) \emph{and} is
        trending down; execute only when every remaining consumer has
        migrated or formally accepted the cutoff. Listing~9.5 automates the
        verdict.
  \item \textbf{Compare JSON and Protocol Buffers for schema evolution.}
        JSON tolerance is client discipline --- the format cannot stop a
        strict deserializer from rejecting unknown members, so contracts must
        mandate tolerant readers. Protobuf tolerance is a wire property:
        unknown tags are skipped by the parser and retained on reserialize,
        and \texttt{reserved} makes name/tag reuse impossible
        \cite{protobufdocs}. JSON evolves by convention; Protobuf by
        construction (Chapter~4).
  \item \textbf{How many versions should you support concurrently?}
        As few as the consumer base permits: two (current + previous) is the
        common external policy; three for slow enterprise bases; date-based
        schemes cap by age instead of count. Each added version pays adapter,
        docs, CI, and pairwise-testing costs that grow superlinearly
        (Equation~\ref{eq:ch9-cost}).
  \item \textbf{How do internal and external API versioning policies differ?}
        Internal consumers are enumerable and often deployable by your own
        org, so dual-serving windows of weeks and aggressive retirement are
        feasible. External consumers are unenumerable and uncontractable at
        scale, so timelines stretch to quarters/years, communication needs
        full campaign machinery, and brownouts replace direct outreach for
        the long tail.
  \item \textbf{What are brownouts and when are they justified?}
        Scheduled, announced windows in which a deprecated surface returns
        errors, used late in a migration campaign to surface silent
        dependencies while rollback is still safe. Justified only after
        direct outreach, with clear notice, and never near the sunset date
        itself.
  \item \textbf{How do you handle a breaking change to a stored representation
        (e.g., cached payloads, event log) while versioning the API?}
        Decouple wire version from storage version: persist in the newest
        schema with upcasting on read, serve old wire versions via adapters,
        and never let a retired wire contract dictate the storage format.
        This is the DDIA dual of expand-and-contract applied to data at rest
        \cite{kleppmann2017ddia}.
\end{enumerate}

\subsection*{Staff (Q28--Q34)}
\begin{enumerate}[label=\textbf{Q\arabic*.}, leftmargin=1.6cm, start=28]
  \item \textbf{Design a version policy for an organization with 200 internal
        APIs and 15 external ones.}
        Federated guardrails: a single published policy defining breaking
        changes, mandatory CI diff gates on every spec repo, internal APIs
        capped at two versions with 30--90 day dual-serving windows enforced
        by platform tooling, external APIs on contractual 12-month sunsets
        with a consumer registry, and an exceptions board owning the
        security-break clause. Uniform telemetry format so retirement
        decisions are mechanical everywhere.
  \item \textbf{When is versioning the wrong tool?}
        When the change can be expressed additively, when consumers are few
        and coordinated (migrate in place), or when the real problem is
        missing tolerant-reader discipline. Every version you can avoid
        minting saves its full carrying cost forever.
  \item \textbf{A principal argues for ``versionless'' APIs. Evaluate.}
        Versionless means maximal additive discipline plus guaranteed
        tolerant readers --- achievable in controlled ecosystems (internal,
        strong SDK governance) and effectively what Stripe simulates
        externally. It fails where you cannot govern client parsing behavior;
        there, explicit versions are the honest acknowledgment of
        uncontrolled consumers.
  \item \textbf{How do you price the cost of a version to justify retirement
        to leadership?}
        Quantify Equation~\ref{eq:ch9-cost} per version: adapter maintenance
        hours, CI lane minutes, documentation drift incidents, on-call
        context-switch cost, and the quadratic test matrix; then put the
        telemetry (consumers remaining, their traffic share) beside it.
        Retirement business cases write themselves when both numbers are
        honest.
  \item \textbf{Design the deprecation-communication architecture for a
        platform with 50k API consumers.}
        Layered: wire-level (\texttt{Deprecation}/\texttt{Sunset}/\texttt{Link}
        headers on every response), artifact-level (changelog, migration
        guides, versioned docs), channel-level (portal banners, newsletter,
        status page), and direct (registry-driven outreach to identified
        top consumers, brownouts for the unresponsive tail). Each layer
        catches the consumers the previous one missed.
  \item \textbf{Your gateway must route five versions of fifty endpoints at
        100k RPS. What breaks first, and how do you design for it?}
        Per-version route tables and adapter latency compound; the answer is
        pushing version dispatch to cheap, cacheable rules (URI prefix
        routing at the edge), keeping adapters in-process and allocation-free
        on hot paths, and making version count a managed budget precisely
        because routing and translation cost scales with $n$.
  \item \textbf{Write the policy for the security-exception clause.}
        Scope (severity-defined), authority (who may invoke), consumer
        notification (maximum feasible notice, out-of-band for identified
        consumers), compatibility duty of care (provide patched old-version
        endpoint where feasible), and post-incident retrospective updating
        the contract. A break invoked without this pre-published framework
        is indistinguishable from an accident.
\end{enumerate}

%% ----------------------------------------------------------------------------
\section{External Bibliography \& Further Reading}
\label{sec:ch9-bibliography}

\begin{itemize}
  \item RFC~8594, \emph{The Sunset HTTP Header Field} (and its original
        \texttt{Deprecation} definition) --- the wire-level deprecation
        vocabulary~\cite{rfc8594}.
  \item RFC~9745, \emph{The Deprecation HTTP Header Field} (2025) ---
        revises \texttt{Deprecation} to the \texttt{@}\emph{<unix-timestamp>}
        form; read alongside RFC~8594.
  \item RFC~9110, \emph{HTTP Semantics} --- \texttt{Accept}, content
        negotiation, \texttt{Vary}, and cache-key behavior underpinning every
        versioning placement decision~\cite{rfc9110}.
  \item OpenAPI Specification 3.1 --- the contract description format our CI
        gate diffs~\cite{openapi31}.
  \item Stripe, \emph{API versioning documentation} and the engineering blog
        post \emph{APIs as infrastructure: future-proofing Stripe with
        versioning} --- the canonical date-based versioning write-up.
  \item J.~Gehl, \emph{API Design Patterns} --- versioning and
        backward-compatibility chapters~\cite{gehl2021apidesign}.
  \item M.~Kleppmann, \emph{Designing Data-Intensive Applications}, ch.~4 ---
        schema evolution, rolling upgrades, and the theoretical basis of
        expand-and-contract~\cite{kleppmann2017ddia}.
  \item Google API Improvement Proposals AIP-180 (\emph{Backwards
        compatibility}) and AIP-185 (\emph{API versioning}) --- Google's
        codified policy.
  \item Microsoft REST API Guidelines, \emph{Versioning} section --- a
        second large-scale policy to compare against Google's.
  \item Protocol Buffers documentation, \emph{Updating a Message Type} ---
        tag reservation, unknown fields, and enum open/closed
        semantics~\cite{protobufdocs}.
  \item The \texttt{oasdiff} project (Tufin) --- production-grade OpenAPI
        diffing and breaking-change detection that Listing~9.2 miniaturizes.
  \item FastAPI documentation --- routers, middleware, and
        \texttt{TestClient} used throughout the listings~\cite{fastapidocs}.
\end{itemize}

%% ----------------------------------------------------------------------------
\section{Capstone Projects}
\label{sec:ch9-capstones}

\subsection*{Capstone 9.1 --- Add v2 to an Existing API with an Adapter Layer \hfill [junior]}
\textbf{Objective:} Take any single-version CRUD API (your Chapter~6 project
qualifies) and evolve one resource through a breaking redesign without
breaking v1 consumers.

\textbf{Inputs/Datasets:} The existing API; the synthetic Northwind Robotics
catalog fixture from Listing~9.1 as the domain data.

\textbf{Steps:} (1)~Define the v2 contract for the resource (at least one
rename and one retype). (2)~Implement the v2 model as the service core.
(3)~Write a pure \texttt{to\_v1} adapter and mount both surfaces.
(4)~Add \texttt{Deprecation}, \texttt{Sunset}, and \texttt{Link} middleware
on v1. (5)~Write adapter-fidelity tests proving v1 values derive from v2.

\textbf{Acceptance Criteria:}
\begin{itemize}
  \item[$\square$] v1 and v2 respond correctly for every seeded record.
  \item[$\square$] v1 responses carry all three lifecycle headers; v2 carries
        none.
  \item[$\square$] Adapter tests pass deterministically offline.
  \item[$\square$] No business logic is duplicated between surfaces.
\end{itemize}

\textbf{Stretch Goals:} Add a third, experimental version behind an opt-in
header; measure adapter overhead with a fixed-iteration timing loop.

\subsection*{Capstone 9.2 --- Media-Type Version Negotiation Middleware \hfill [mid]}
\textbf{Objective:} Generalize Listing~9.3 into reusable FastAPI middleware
that version-negotiates \emph{every} endpoint of an API by vendor media type.

\textbf{Inputs/Datasets:} Listing~9.3; the dual-serving app of Listing~9.1
refactored to a single \texttt{/catalog} URI.

\textbf{Steps:} (1)~Implement the negotiator as middleware that resolves the
version and stashes it in request state. (2)~Route representations through an
adapter registry keyed by version. (3)~Emit 406 with a machine-readable
supported-versions body on failure. (4)~Emit \texttt{Vary: Accept} on all
versioned responses. (5)~Test q-value precedence, unsupported versions, and
default selection.

\textbf{Acceptance Criteria:}
\begin{itemize}
  \item[$\square$] All negotiation outcomes match a published decision table.
  \item[$\square$] 406 bodies enumerate supported \texttt{version=N} values.
  \item[$\square$] \texttt{Vary: Accept} present on every 200.
  \item[$\square$] Full pytest suite passes offline.
\end{itemize}

\textbf{Stretch Goals:} Content-negotiate a second wire format
(\texttt{+json} vs.\ \texttt{+cbor}-style suffix) orthogonally to version.

\subsection*{Capstone 9.3 --- Production Breaking-Change CI Gate \hfill [mid]}
\textbf{Objective:} Harden Listing~9.2 into a gate you would actually run on
pull requests.

\textbf{Inputs/Datasets:} Listing~9.2; the two catalog fixture specs; three
additional spec pairs you author (purely additive, parameter-level breaking,
\texttt{\$ref}-based).

\textbf{Steps:} (1)~Add \texttt{\$ref} resolution for local references.
(2)~Diff path/query \texttt{parameters} (new required or removed =
breaking). (3)~Classify enum growth as a \emph{warning} tier between breaking
and clean. (4)~Emit a machine-readable JSON report alongside the text one.
(5)~Add a \texttt{--allow-breaking-with=VERSION\_BUMP} escape hatch that
passes only when \texttt{info.version}'s major component increased.

\textbf{Acceptance Criteria:}
\begin{itemize}
  \item[$\square$] Exit code 1 exactly when breaking changes exist without a
        major bump.
  \item[$\square$] All five fixture pairs produce the expected verdicts.
  \item[$\square$] JSON report round-trips through \texttt{json.loads}.
\end{itemize}

\textbf{Stretch Goals:} Detect semantic-risk heuristics (format changes,
pattern tightening, \texttt{additionalProperties: false} introductions).

\subsection*{Capstone 9.4 --- Deprecation Campaign Plan \& Retirement Analytics \hfill [senior]}
\textbf{Objective:} Design and instrument the complete retirement program for
a deprecated version, from announcement to 410.

\textbf{Inputs/Datasets:} Listing~9.5 and its log fixture, extended by you to
6 months and 12 synthetic consumers with realistic migration curves.

\textbf{Steps:} (1)~Write the one-page deprecation policy (dates, SLAs,
channels, brownout schedule). (2)~Extend the report generator with
per-consumer trend lines and a projected-zero-crossing date via linear
fit. (3)~Produce the consumer contact list ranked by dependency.
(4)~Script the brownout simulator: given the log, estimate which consumers
would notice first. (5)~Define the go/no-go checklist for executing sunset.

\textbf{Acceptance Criteria:}
\begin{itemize}
  \item[$\square$] Policy document covers timeline, SLAs, channels, and
        exceptions.
  \item[$\square$] Report projects a retirement date and names remaining
        consumers deterministically.
  \item[$\square$] Go/no-go checklist is mechanical (every item checkable
        from telemetry).
\end{itemize}

\textbf{Stretch Goals:} Add a ``what-if'' mode showing how the projected date
shifts if the top two consumers migrate next month.

\subsection*{Capstone 9.5 --- Version Governance at Platform Scale \hfill [staff]}
\textbf{Objective:} Design the versioning and retirement governance system
for a fictitious platform org running 40 internal and 6 external APIs.

\textbf{Inputs/Datasets:} A synthetic portfolio registry you author (YAML:
per-API versions, consumer counts, traffic shares, owners); this chapter's
cost model (Equation~\ref{eq:ch9-cost}).

\textbf{Steps:} (1)~Codify the policy: breaking-change definition, allowed
strategies per audience, sunset floors, exception clauses. (2)~Build a
registry linter that flags policy violations (version in URI \emph{and}
header; third live version without waiver; deprecated version lacking a
sunset date). (3)~Compute the portfolio's carrying cost under your model and
identify the five most expensive versions. (4)~Write the RFC-style decision
record proposing your org-wide strategy mix. (5)~Define the telemetry schema
every API must emit.

\textbf{Acceptance Criteria:}
\begin{itemize}
  \item[$\square$] Policy document is complete enough to lint mechanically.
  \item[$\square$] Linter flags at least the four violation classes above on
        a seeded registry.
  \item[$\square$] Cost report ranks versions and is reproducible offline.
  \item[$\square$] Decision record defends the strategy mix against at least
        two alternatives.
\end{itemize}

\textbf{Stretch Goals:} Simulate ten years of portfolio growth and show when
the quadratic term of Equation~\ref{eq:ch9-cost} forces policy change.

%% ----------------------------------------------------------------------------
\section{Python Dependencies Appendix}
\label{sec:ch9-deps}

All code runs offline after a one-time install. From PowerShell:

\begin{minted}{text}
python -m venv .venv
.\.venv\Scripts\Activate.ps1
pip install -r requirements.txt
\end{minted}

\begin{minted}{text}
fastapi>=0.110
pydantic>=2.7
pyyaml>=6.0
jsonschema>=4.21
httpx>=0.27
pytest>=8
\end{minted}

\subsection{fastapi ($\geq$ 0.110)}
\textbf{Description:} The ASGI web framework used for every served API in
this chapter. \textbf{Why included:} routing, Pydantic integration,
middleware (deprecation headers), and \texttt{TestClient} for offline
demonstrations~\cite{fastapidocs}. \textbf{Alternatives:} Flask + manual
validation; Starlette directly (FastAPI's base). \textbf{Installation:}
\texttt{pip install "fastapi>=0.110"}.

\subsection{pydantic ($\geq$ 2.7)}
\textbf{Description:} Type-driven data validation and serialization library.
\textbf{Why included:} defines the v1/v2 wire models and the domain model in
Listing~9.1; its strict-by-default validation is the producer-side
counterpart to the chapter's consumer tolerance rules. \textbf{Alternatives:}
\texttt{dataclasses} + manual checks; attrs; marshmallow.
\textbf{Installation:} \texttt{pip install "pydantic>=2.7"}.

\subsection{pyyaml ($\geq$ 6.0)}
\textbf{Description:} YAML parser/emitter. \textbf{Why included:} the
breaking-change gate (Listing~9.2) reads OpenAPI descriptions authored in
YAML, the format real API teams use~\cite{openapi31}. \textbf{Alternatives:}
ruamel.yaml (round-trip fidelity); parsing JSON-flavored OpenAPI with stdlib
\texttt{json}. \textbf{Installation:} \texttt{pip install "pyyaml>=6.0"}.

\subsection{jsonschema ($\geq$ 4.21)}
\textbf{Description:} JSON Schema validation engine. \textbf{Why included:}
exercise and capstone work validating consumer fixtures against published
schemas (E9.2, Capstone~9.3's semantic checks); not required by the printed
listings. \textbf{Alternatives:} manual validation; Pydantic's schema export.
\textbf{Installation:} \texttt{pip install "jsonschema>=4.21"}.

\subsection{httpx ($\geq$ 0.27)}
\textbf{Description:} Sync/async HTTP client. \textbf{Why included:} the
transport underneath FastAPI's \texttt{TestClient} --- required for
Listings~9.1, 9.3, and 9.4 to run offline. \textbf{Alternatives:} requests
(synchronous, not ASGI-native); urllib (stdlib, low-level).
\textbf{Installation:} \texttt{pip install "httpx>=0.27"}.

\subsection{pytest ($\geq$ 8)}
\textbf{Description:} Test framework. \textbf{Why included:} runs the
consumer-tolerance suite (Listing~9.4) and every capstone's acceptance
checks. \textbf{Alternatives:} unittest (stdlib); nose2.
\textbf{Installation:} \texttt{pip install "pytest>=8"}.

%% ----------------------------------------------------------------------------
\section{Chapter Summary \& Key Takeaways}
\label{sec:ch9-summary}

\begin{itemize}
  \item A breaking change is defined by its victims: any change that
        invalidates a formerly successful client interaction. Renames, type
        changes, semantic drift, validation tightening, and required-field
        additions are all breaking; additions and loosening are not ---
        except enum growth against closed consumer parsers.
  \item Postel's law, applied maturely, means strict about structure and
        safety, tolerant about extensibility. Tolerant readers make additive
        evolution free; blind tolerance fossilizes bugs into contracts.
  \item Four versioning placements --- URI path, query parameter, custom
        header, media type --- trade cache correctness, gateway simplicity,
        debuggability, and REST purity differently. Stripe (dates), GitHub
        (header + media type), and Google (URI) each chose rationally for
        their consumer base; none is universally right.
  \item Expand-and-contract is the only honest zero-downtime migration
        shape: add beside, migrate with telemetry, contract on evidence.
        Dual-serving keeps $n$ surfaces affordable by putting one domain
        model behind pure, tested adapters.
  \item Deprecation is engineered, not announced: \texttt{Deprecation}/
        \texttt{Sunset}/\texttt{Link} headers (RFC~8594, revised by
        RFC~9745), published timelines, stability SLAs, SemVer changelogs,
        and layered communication~\cite{rfc8594}.
  \item Schema discipline is additive-only with reserved names and open
        enums; JSON tolerates by convention, Protobuf by construction
        \cite{protobufdocs}.
  \item Compatibility belongs in CI: an OpenAPI diff gate fails the build on
        breaking changes before consumers do.
  \item Every live version carries recurring, superlinear cost. Retirement
        is the goal state; telemetry --- not the calendar --- grants
        permission to reach it.
\end{itemize}

%% ----------------------------------------------------------------------------
\section{Quick-Reference Card}
\label{sec:ch9-quickref}

\subsection*{Breaking-Change Checklist}
Breaking: removed/renamed field $\bullet$ type or semantics change $\bullet$
tightened validation $\bullet$ new required request field $\bullet$ removed
endpoint/operation $\bullet$ changed URL $\bullet$ changed error
codes/shape $\bullet$ enum value \emph{removed} $\bullet$ enum value
\emph{added to a documented-closed enum}. Non-breaking: added endpoint
$\bullet$ added optional request field $\bullet$ added response field
(tolerant readers) $\bullet$ loosened validation $\bullet$ added enum value
to a documented-open enum.

\subsection*{Strategy Matrix (Condensed)}
URI path --- best ops ergonomics, worst REST purity $\bullet$ query param ---
URI-like, weaker tooling $\bullet$ custom header --- pure URIs, fragile
caches, needs \texttt{Vary} $\bullet$ media type --- most REST-orthodox,
highest consumer complexity, needs \texttt{Vary: Accept}. Pick one signal;
never two.

\subsection*{Deprecation-Header Reference}
\texttt{Deprecation: @}\emph{<unix-ts>} (RFC~9745; RFC~8594 used HTTP-date)
--- when deprecation began. \texttt{Sunset: }\emph{<HTTP-date>} (RFC~8594)
--- when the surface dies. \texttt{Link: <...>; rel="deprecation"} --- policy
document. \texttt{rel="successor-version"} --- where to migrate. Lifecycle:
introduced $\to$ current $\to$ deprecated $\to$ sunset $\to$ retired (410).

\section{Standardized Evidence Addendum}
\subsection{Benchmark Snapshot Template}
\begin{table}[h]
  \centering
  \small
  \begin{tabularx}{\textwidth}{l c c c c}
    \toprule
    \textbf{Config} & \textbf{p50 latency} & \textbf{p99 latency} & \textbf{Error rate} & \textbf{Throughput} \\
    \midrule
    Baseline & -- & -- & -- & 1.00x \\
    Candidate A & -- & -- & -- & -- \\
    Candidate B & -- & -- & -- & -- \\
    \bottomrule
  \end{tabularx}
  \caption{Minimum benchmark table to keep per chapter.}
\end{table}

\subsection{Request Trace Template}
\begin{itemize}
  \item \textbf{Request:} one representative API call for this chapter's topic.
  \item \textbf{Before:} behavior (headers, payload, latency, failure mode) with the naive approach.
  \item \textbf{After:} behavior with the technique in this chapter.
  \item \textbf{Result:} what changed in correctness, resilience, latency, and operability.
\end{itemize}

\subsection{Cost and Latency Budget Snapshot}
\begin{table}[h]
  \centering
  \small
  \begin{tabularx}{\textwidth}{l c c}
    \toprule
    \textbf{Stage} & \textbf{Latency budget} & \textbf{Cost driver} \\
    \midrule
    TLS / connection setup & -- & handshake round trips, keep-alive hit rate \\
    Request validation & -- & schema complexity, CPU per request \\
    Business logic and I/O & -- & downstream fan-out, query cost \\
    Serialization and egress & -- & payload size, compression ratio \\
    \bottomrule
  \end{tabularx}
  \caption{Operational budget template for this chapter's technique.}
\end{table}

\section{Configuration Preset Template}
\begin{minted}{text}
# api_policy.yaml
policy_version: "v1.0.0"
component: "<chapter_topic>"
timeout_ms: 0
retry_budget: 0
fallback_strategy: "fail_fast"
rollback_trigger: "error_budget_burn_gt_threshold"
\end{minted}

\section{CI Quality Gate Specification}
\begin{itemize}
  \item contract tests green against the pinned OpenAPI/AsyncAPI/Protobuf spec,
  \item no breaking schema changes without a version bump,
  \item error responses conform to RFC 9457 problem-details shape,
  \item benchmark deltas within approved regression bounds (latency and throughput),
  \item policy version and rollback plan present in release metadata.
\end{itemize}

\section{Incident Runbook Template}
\begin{enumerate}
  \item Detect regression from SLO burn-rate alerts or consumer reports.
  \item Scope impacted endpoints, versions, and consumer segments.
  \item Contain impact (rollback, feature flag, rate-limit tightening, or traffic shift).
  \item Diagnose with traces, structured logs, and contract diffs.
  \item Recover with validated fix and verified rollout procedure.
  \item Prevent recurrence via new regression tests and CI gates.
\end{enumerate}

\section{Cross-Chapter Decision Card}
\begin{table}[h]
  \centering
  \small
  \begin{tabularx}{\textwidth}{l X X}
    \toprule
    \textbf{Dimension} & \textbf{Choose this approach when} & \textbf{Prefer an alternative when} \\
    \midrule
    Use case & -- & -- \\
    Latency budget & -- & -- \\
    Operational cost & -- & -- \\
    Team maturity & -- & -- \\
    \bottomrule
  \end{tabularx}
  \caption{Decision card to complete per chapter.}
\end{table}

\section{ROI Model and Build-vs-Buy}
\begin{itemize}
  \item \textbf{Build when:} the capability is differentiating, compliance forbids vendors, or scale makes vendor pricing irrational.
  \item \textbf{Buy when:} the capability is commodity (auth, gateways, observability), time-to-market dominates, or staffing is thin.
  \item \textbf{Hidden costs to model:} on-call load, upgrade cadence, expertise ramp, data egress fees.
\end{itemize}

\section{Disaster Recovery Notes (RPO/RTO)}
\begin{itemize}
  \item Define recovery point objective (RPO): maximum acceptable data loss window.
  \item Define recovery time objective (RTO): maximum acceptable restoration time.
  \item Test restores regularly; an untested backup is not a backup.
\end{itemize}

